% ============================================================
%  Harald Høffding, "Erkendelsesteori og Livsopfattelse"
%  "Epistemology and the Conception of Life"
%  Copenhagen: Det Kongelige Danske Videnskabernes Selskab, 1925.
%  (Source text: Lindhardt og Ringhof e-book edition, 2020.)
%  TRANSLATION (English)
%
%  Høffding's running paragraph numbers (1, 2, 3 ...) within each
%  chapter are preserved as bold run-in numerals. His scholarly
%  notes are rendered as footnotes; cited work-titles are kept in
%  the original language (italic) per normal bibliographic practice,
%  with an English gloss where helpful.
% ============================================================
\documentclass[12pt]{book}
\usepackage[utf8]{inputenc}
\usepackage[T1]{fontenc}
\usepackage[english]{babel}
\usepackage[protrusion=true,expansion=true]{microtype}
\usepackage[margin=1.3in]{geometry}
\usepackage{setspace}
\usepackage{amsmath}
\usepackage{libertinus}
\usepackage{libertinust1math}
\usepackage{fancyhdr}
\usepackage[colorlinks=true,linkcolor=black,urlcolor=black,
  pdftitle={Epistemology and the Conception of Life},
  pdfauthor={Harald Høffding}]{hyperref}
\renewcommand{\thechapter}{\Roman{chapter}}
\setlength{\headheight}{15pt}
\pagestyle{fancy}
\fancyhf{}
\fancyhead[LE,RO]{\thepage}
\fancyhead[RE]{\textit{Epistemology and the Conception of Life}}
\fancyhead[LO]{\textit{\leftmark}}
\setlength{\parindent}{1.5em}
\onehalfspacing

% run-in section number (Høffding's paragraph numbering)
\newcommand{\hnum}[1]{\textbf{#1.}\enspace}

\title{Epistemology and the Conception of Life}
\author{Harald Høffding}
\date{Copenhagen, 1925\\[6pt]
  \small Translated from the Danish}

\begin{document}
\maketitle

\chapter{Epistemology}
\label{ch:1}

\hnum{1} According to the conception of philosophy that has gained the upper hand since
Kant's time, philosophy is an investigation of the presuppositions of human spiritual life.
Human spiritual life lays itself before the day in science, art, morality, and religion. In all
these domains there lie presuppositions at the foundation which need not appear clearly and
with full consciousness to those who live and work in them. Philosophy can produce neither
special science, art, morality, nor religion, but it can perhaps show that in all these domains
one works, consciously or unconsciously, with thoughts that it may have its interest to
investigate, especially if it can be shown that in these thoughts there betrays itself a
structure peculiar to human spiritual life. If it succeeds in drawing forth such thoughts, the
question becomes whether we have in them revelations of the innermost essence of existence, or
whether they merely stand as forms which human striving in the various domains is, after all,
referred to employing. Already in Plato there appears the tragedy that presents itself at every
attempt to demonstrate an absolute foundation for the work of spiritual life and its fortunes.
Again and again Plato attempts to derive all cognition and all existence from a highest idea
that should not itself need any grounding. In his books ``on the State'' one can follow his
striving to reach up to what for him would have to be the absolute standpoint. Again and again
he is thrown back in his approaches to reach this position, and he in fact admits that only by
the help of poetic analogies can he find expression for the goal that stood before
him.\footnote{Cf.\ \textit{Platon's Bøger om Staten. Analyse og Karakteristik} [Plato's Books on
the State: Analysis and Characterization] (Vidensk.\ Selsk.\ filos.\ Meddelelser, 1924), pp.\
87--110.} Two thousand years later Kant established, with full consciousness, what Plato so
reluctantly would admit.

Although there exist significant philosophical attempts to investigate the presuppositions of
art, morality, and religion, it was yet the investigation of the foundation of science that
came to play the largest role in philosophy. Epistemology has more and more become the central
discipline of philosophy; indeed, one can even trace a tendency to identify philosophy and
epistemology, so that it becomes a great question what position the disciplines that treat the
other domains of spiritual life can attain within philosophy. These other domains do, however,
also raise their problems, which cannot without further ado be treated merely from the
standpoints of epistemology. As thinkers in the face of whom such a question must be raised, one
must especially name Francis Bradley (\textit{Appearance and Reality}, 1893), for whom the
essence of philosophy is properly exhausted with the investigation of what measure we have for
distinguishing between reality and phenomenon, and Hermann Cohen (\textit{Logik der reinen
Erkenntnis}, 1902), founder of the Marburg School.\footnote{Cf.\ on these thinkers: \textit{Den
nyere Filosofis Historie} [History of Modern Philosophy], 3rd ed., IV, pp.\ 210\,f.\ and
217\,f.}

But there is yet another question that in the most recent time presses forward more and more.
From that deep root of the mind from which not merely science, but art, morality, and religion
spring, there proceed also other psychic formations, involuntary spiritual crystallizations, in
which men get expression for what stirs in their minds, and for their real or supposed
experiences. These are views of life and of the world, more or less formed into peculiar
wholes, which cannot without further ado be reckoned under one of the four principal forms of
spiritual life named in the foregoing.

Just as little as philosophy can produce special science, art, morality, and religion, but must
see its task in investigating the presuppositions of the whole-formations that appear in these
domains, just as little is it able to produce the individual whole-formations that appear as
views of life or of the world; rather it must provisionally take them as something given,
something at hand, syntheses factually produced, and then afterward seek to ascertain by what
ways and by what stages they have come to be, and what conditions their continued subsistence
requires. It must then resolve what has involuntarily been built up, and discuss its
significance for human spiritual life as a whole. ---

There are many ways that can lead to philosophy, and the individual philosopher's position
toward its various problems will for an essential part be determined by the way along which he
first came to philosophize. For my part, the interest was from the first bound up with
psychological, ethical, and religious questions. I therefore worked out the three philosophical
disciplines corresponding to these questions (1876--1901). But precisely through the treatment
of these three problems I was led to see the fundamental position of the problem of cognition in
relation to them. --- Psychology approaches this problem by the question of the conditions for
the validity of our thinking. Both the valid and the invalid thinking are, after all, in fact
thinking, and can so far be investigated psychologically. By the question of the difference
between validity and invalidity the transition from psychology to epistemology takes place. But
besides this, psychological research itself, like all other research, rests upon presuppositions
which it is epistemology's business to discuss. --- In the attempt to give a grounding of
ethical judgments it shows itself that, when one proceeds strictly logically, there are
different typical starting-points, and that an epistemological discussion must set certain
limits to the scientific character of ethics. --- In the religious domain the question becomes
whether religion in any of its forms can develop thoughts that can be laid as a foundation in
the interpretation of existence, as the latter is known through experience and is accessible to
science.

In \textit{Den menneskelige Tanke} [The Human Thought] (1910) I then go directly into
epistemological investigation, which I have sought to continue, as regards a few important
concepts, in three treatises (on Totality, Relation, and Analogy).

While the more concrete disciplines all presuppose questions that it becomes epistemology's
business to discuss, one cannot conversely derive the subjects of the concrete disciplines from
epistemological standpoints. And this holds not merely for disciplines that stand in close
relation to philosophy, such as psychology, ethics, and the philosophy of religion. It holds for
every special science. Epistemology presupposes that there is science in the world. Plato
already saw this when he found the task of philosophy in analyzing the presuppositions on which
mathematics, the only special science he knew, rests.\footnote{Cf.\ \textit{Platon's Bøger om
Staten}, pp.\ 98; 107\,f.} Epistemology has so far the character of purely logical analysis. But
if it is asked how the peculiar spiritual phenomenon we call science has in fact come to be,
there is thereby set a task for psychology on the basis of the history of science. In the
fundamental concepts (categories) with which scientific research works, it shows itself how the
human spirit's need and capacity are formed in special ways in reciprocal action with the
experiences present at each given time.\footnote{Cf.\ my account of the doctrine of categories
in \textit{Den menneskelige Tanke}, pp.\ 135--245.} Naturally this psychological-historical
method can and must itself again become the object of epistemological investigation; the problem
of cognition arises continually anew, for the good reason that every attempt to treat it,
whatever method one then employs, rests upon certain presuppositions.

While psychological-historical investigation within epistemology takes a preparatory and
subordinate place in relation to purely logical analysis, it is otherwise in the more concrete
domains on which philosophical thinking can work. In the face of art, morality, and religion,
understanding must be sought by the psychological-historical way of how such wholes --- as every
work of art, every practical morality, every living religion is --- have been able to come to
be, and how far the conditions that have here been decisive can be assumed to be present also in
the future. And in a similar way philosophy stands in the face of individual views of life.
These too make use of thoughts and rest upon certain presuppositions. Even where such a view of
life is, or supposes itself to be, independent of all science, the inexorable law of life, and
especially the inner structure to which all spiritual life bears witness, will yet demand a
certain coherence, not merely among the thoughts that are woven into the view of life, but also
between these thoughts and those the individual otherwise makes use of in his struggle for
existence. In every serious view of life there will be traceable a need to find a collected
expression for the demands, the tasks, and the fortunes that the individual's life carries with
it.

It will be a particular task of this treatise to investigate the character of such individual
views of life and their relation to speculation (``metaphysics'') on the one side, and to
religion (especially positive, handed-down religion) on the other. Individual views of life are
often, like art, morality, and religion, regarded by philosophers and natural researchers merely
as ``anthropomorphisms,'' as spiritual phenomena that from the standpoint of strict thinking have
no interest. Parmenides already distinguishes sharply between the pure identity that thought
seeks and maintains, and ``the opinions of mortals,'' which are due to contingent and shifting
standpoints. But philosophy must investigate whether behind these anthropomorphisms there cannot
be traced an involuntary spiritual labor, a witness to the inner structure that all spiritual
life, even the strictest science, presents.

In my account of the history of modern philosophy I have often investigated the philosophers'
individual views of life in their relation to their systems; in my book on Kierkegaard I have
especially gone into this question, as also in my account of Nietzsche (in \textit{Moderne
Filosofer} [Modern Philosophers]), in my comparison between Kierkegaard and Nietzsche (in the
third volume of my \textit{Mindre Arbejder} [Minor Works]) and between Pascal and Kierkegaard
(\textit{Tilskueren}, 1923). Besides this, in my psychological study \textit{Den store Humor}
[The Great Humour] I have sought to characterize a peculiar type of the way of taking life. ---
In the following investigation I keep to the domains I have thought about most. Therefore I do
not go into the problems of art, however great a personal significance art has otherwise had for
me.

\hnum{2} As I have sought to show in my treatise \textit{Totalitet som Kategori} [Totality as a
Category] (Vidensk.\ Selsk.\ Skr., 1917), the thought of the whole plays a double role in our
cognition. In the first place, all cognition is a striving to connect observations and
representations into a whole. It is a struggle with chaos, with all that is scattered and
incoherent. The point is not merely to connect, but to connect in such a way that the subject of
cognition appears in definite relations, partly between the subject's own elements, partly
between the subject as a whole and other subjects. Here synthesis and relation announce
themselves as the first in the series of categories. To these first categories there attach
themselves continuity and discontinuity, likeness and difference, in mutual conflict.
Discontinuity and difference rouse thought to activity, but the activity then sets out to find
continuity behind the discontinuity, likeness behind the differences. Such conflict already takes
place in ordinary human understanding, the pre- and unscientific consciousness. The three pairs
of concepts named I therefore reckon among the fundamental categories that make themselves felt
in all conscious life and are thus common to the pre- and unscientific and the scientific
cognition. The so-called formal and real categories are more special and more exact
determinations of the fundamental ones.

The possibility of forming complete thought-wholes is not the same in all domains. In logic and
mathematics more perfect thought-wholes can be formed than in domains where ever-new observation
is necessary. In a strict logical inference, where the relation between premises and conclusion
stands quite clear, we have an example of a rational whole, especially when the inference can be
reversed, so that from the conclusion one can derive the premises themselves. Every link in such
a thought-whole is conditioned by the other links and thereby by the whole itself, and
conversely. Such a transparent whole has at all times stood as the great ideal of cognition,
whether one has believed that this ideal could be realized or not. And it is an ideal that hangs
closely together with the first categories in the fundamental group --- synthesis and relation.
The closest connection between thought-elements is entered upon through the fully clarified
relations. It was in order to make an approximation to this ideal possible that Plato, through a
purification and raising-to-a-higher-power of the relations of likeness that observations can
present, reached the strict concept of identity which he designated with the word ``idea.'' And
an analogous process of purification has in recent times, not without the influence of Plato,
taken place with the concepts number, time, place, and degree, in order to make them fit to enter
the service of the new natural science.\footnote{Cf.\ \textit{Den menneskelige Tanke}, pp.\
181--201.}

Strict science has developed out of ``anthropomorphisms.'' Ordinary consciousness sets out from
the assumption that the world is as it presents itself to us, with great wealth of qualities and
with a multitude of wholes that stand more or less sharply over against one another. Here the
whole appears, not for science as an ideal that is to be reached, but as a problem that is to be
solved. It is science's task to describe such wholes and then to find the conditions for their
coming-to-be and continued subsistence. This is the second side in science's relation to the
concept of the whole.

But science itself does not constitute a strict whole. It divides into different parts according
to the different domains it works on, and these parts appear, on account of the nature of the
matters, with a highly different degree of strict scientific character. The one science can
thereby acquire an irrational, ``anthropomorphic'' character in relation to the other. One will be
able to arrange the sciences in a series such that each single link stands with greater strictness
in relation to the following, while in relation to its predecessor it acquires a more imperfect
character --- strictness or perfection being constantly measured by the approximation to the
carrying-through of the principle of identity. Thus geometry stands in relation to arithmetic,
physics in relation to mathematics, biology in relation to physics, psychology and the other
sciences of spirit in relation to biology.\footnote{Cf.\ \textit{Begrebet Analogi} [The Concept
of Analogy], pp.\ 80--88.} The relation between the rational and the irrational
(anthropomorphic) extends right up into the series of the sciences; it appears not merely in the
relation between science on the one side and the current representations on the other. In the
relation between the sciences among themselves, rationality can yet be carried through in the
sense of rational analogy. Thus when the geometrical forms are apprehended as symbols of what
only arithmetic can express with strict exactness. Real science rests upon analogy between
logical-mathematical series and series of observation. As Kant has shown, the ``application'' of
logical and mathematical forms to subjects at hand consists in a rational analogizing.\footnote{Cf.\
\textit{Begrebet Analogi}, pp.\ 89--109.} Without here again going into what I have discussed in
detail in earlier treatises, I shall merely remark that our criterion of reality in doubtful cases
consists in the possibility of working the subjects at hand together, so that they show themselves
to stand in rational connection. Scientifically we seek to work through such a connection in such a
way that it approaches as nearly as possible to a purely logical or mathematical relation and
becomes transparent like these. At various points in the history of thinking we find even
tendencies to make reality into pure logic or pure mathematics. Émile Meyerson has recently\footnote{É.\
Meyerson, \textit{Déduction relativiste} (1925), passim.} demonstrated a tendency in this direction
in some of Einstein's adherents, in likeness with what Parmenides, Descartes, and Hegel attempted
from purely philosophical standpoints. ---

Despite everything, the type of all thought-life is the same. We do not, in science, put on a wholly
different way of thinking than the one we employ in our practical, personal life. Science is a form
of spiritual life and bears the characteristic stamp of all spiritual life.

The epistemology developed on the foundation given by Kant, the so-called Criticism, distinguishes
itself on this point from the so-called Pragmatism, which lays the chief weight on the outer or inner
motives for maintaining one's thoughts, without these thoughts' form and character standing in an
essential relation to the nature and structure of the human spirit. Our scientific concepts, the
fundamental concepts too, are indeed to correspond to human need, but human capacity is not to be
cognizable through them. It is, to be sure, only in instinct and genius that need and capacity
accompany each other all the way through. But a constant disharmony between them will, under the
human race's continued life, be impossible. It is now not merely need that can make itself felt
without corresponding capacity; it is also the capacity of thought that can unfold itself without any
real human need being felt to be satisfied thereby. In the first case the heart's hope or fear, its
wish or its aversion, tries to dictate to thinking its results. In the second case the heart's free
life is forced back in the face of real or supposed intellectual conquests. A solution of such
conflicts becomes possible only by the fact that both human need and human capacity are subject to
development. Thus, for example, the need to obtain at any price a solution for all riddles, because
life otherwise could not be endured, can undergo a metamorphosis that is brought about by experience
that what makes life worth living is not least the labor of thought, the use of intellectual
capacity, to whatever goal it then leads. As regards capacity, there has not seldom been a tendency
to believe that one should be able, once and for all, to decide how far it can lead. Even Kant
supposed that the account of human cognition must be capable of being settled once and for all. But
there is no doubt that human capacity for cognition (``reason'') undergoes development, just as well
as human need.

That all men (researchers included) have not at all times had the same ``reason'' is shown by the
history of thinking. A good example is to be had in the fact that the discoverers of the proposition
of the conservation of energy, toward the middle of the nineteenth century, all each for himself
declared that this proposition was really self-evident, since there is nothing that can become
nothing, just as little as anything can come to be out of nothing. But if one asks whether this
self-evidence has been acknowledged at all times, history shows that ancient physics was built up on
the opposite assumption, and that it was first the thinkers of the seventeenth century who maintained
that there could be no more in the effect than in the cause, an assumption that the carrying-through
of the proposition of the conservation of energy establishes in detail.\footnote{Cf.\ \textit{Ledende
Tanker i det nittende Aarhundrede} [Leading Thoughts in the Nineteenth Century], pp.\ 27\,f.}

When now the ``self-evidence'' varies from time to time, that means that ``reason'' varies. The sound
human understanding, from which science has sprung, follows slowly, sometimes very slowly, after
science in its development, and gradually takes up into itself, as self-evident, the most accessible
scientific results. There are those who regard it as self-evident that one telegraphs through a wire,
but do not understand that there can also be wireless telegraphy.

Émile Meyerson gives an interesting example of the variation of ``reason,'' drawn from the development
of the apprehension of space. While in antiquity one was inclined to suppose that the direction ``up
and down'' took a special position in relation to the other directions in space, the new natural
science maintained that this special position did not hold, since space was uniform in all
directions. Later one long regarded it as a matter of course that space had only three dimensions.
But there has then developed a non-Euclidean geometry, which plays an essential role in the newest
discoveries of physics.\footnote{\textit{Déduction relativiste}, pp.\ 317\,f.} Meyerson supposes,
however, that outside geometrical reason no clear examples of changes of ``reason'' can be
demonstrated. Against this, I think, speaks the above-cited example with the proposition of the
conservation of energy. And, on the whole, if one apprehends the fundamental concepts of cognition
(the categories) as thoughts that are drawn forth by study of the work of cognition at its various
stages --- a conception I have set forth in \textit{Den menneskelige Tanke}, and in which I believe
Meyerson agrees with me\footnote{\textit{Déduction relativiste}, pp.\ 170; 308\,f.} --- then there is
nothing to prevent categories from being able to die out, just as they can come to be. They die away
when the work of cognition, in the face of new observations that demand interpretation, is obliged to
strike out on new ways. Such a dying, if not dead, category is the concept of substance, from the
very moment when Leibniz declared that we know substances only as acting, and that everything that
acts is substance.\footnote{\textit{Den menneskelige Tanke}, p.\ 217.} The point is then to find the
law for the activity appearing in our experience; through this law we have the essence of the acting
thing expressed. This holds whether there is talk of spiritual or of material substances. Already in
this it lies that it is not merely in the natural-scientific, but also in the spiritual-scientific
domain, that ``reason'' shows itself to be in development. Not least does it show itself in the manner
in which religious and social questions are nowadays treated in serious discussion. One does not set
out from certain thoughts fixed once and for all, but investigates the connection of religious and
social phenomena with man's nature and with the historical conditions, and is clear that when these
are changed, religion and social culture too must be changed. Our ``reason'' is now, in contrast to
the ``reason'' of earlier times, saturated with psychological and historical insight. Already a
comparison between Feuerbach's and Voltaire's discussions of religious questions shows this.

After this conception of the fundamental concepts that make themselves felt in the human striving
after cognition, a clear light falls upon Lessing's well-known proposition that the constant striving
after truth is better than the very possession of truth --- that the hunt is better than the
quarry.\footnote{On the understanding of this proposition see \textit{Den nyere Filosofis Historie},
III, 3rd ed., pp.\ 19\,f.} Lessing's grounding is that only by striving, not by possession, are our
powers exercised and extended. And when now, in the history of cognition, it so often shows itself
that the possession one supposed oneself to have reached was illusory, or in any case had significance
only as a starting-point for new work, then the ever-continued striving steps forth so much the more
clearly in its worth. Even if one has come upon a false track, the consistent pursuit of the track can
yet in any case teach that one must strike out on other tracks in order to get forward. Consistent
thinking constantly brings its reward with it. The ideal of cognition is, as we have seen, a rational
whole, won through the mutual determination of all links, and it is an ideal that will constantly set
new tasks, partly of a positive, partly of a negative kind, partly about ways, partly about strayings.
There is, in the end, not the opposition between striving and possession that Lessing set out from.
The flying arrow is at each moment both at rest, in that it is constantly at a definite place, and in
motion, since it is constantly in the act of leaving this place.

\hnum{3} Every scientific proposition corresponds to a certain stage of development of human reason,
and since the sound human understanding cannot as quickly as science reach this stage of development,
every proposition cannot without further ado be made evident to it. Here there is constantly
possibility of doubt and strife. Often a single observer or thinker must for a time stand as guarantor
for a cognition that cannot at first hand be accessible to all. There may then be doubt whether he has
observed or thought correctly. It must be tested --- by other observers or thinkers. And the testing
must take place from presuppositions that, at the given stage of reason's development, stand as
necessary for human cognition. If now these presuppositions become the object of doubt, one must try
to find other presuppositions whose necessity can better be grounded, although they too of course
correspond to a certain stage of development of human reason. There is, after all, no absolute, but
only a hypothetical necessity. And it must here not be forgotten, what Descartes already drew
attention to, that the very act of doubting is a thinking that, like all thinking, rests upon certain
presuppositions. If these presuppositions fall before criticism, then doubt too falls away, just as
well as dogmas fall away when the foundation on which they build falls.

From this consideration it follows that it will not be possible, once and for all, to draw boundaries
between anthropomorphisms and absolutely valid representations. The thinking too that should teach us
something about ``the things in themselves'' is human thinking and rests upon certain presuppositions.
I here leave aside that anthropomorphisms too demand to be explained. How are they at all possible?
Spinoza has rightly inculcated that all our representations, the valid as well as the invalid, arise
with the same necessity\footnote{\textit{Omnes ideæ, tam adæquatæ quam inadæquatæ, eadem necessitate
consequuntur} [All ideas, the adequate as well as the inadequate, follow with the same necessity].
\textit{Ethica} II, 36, Dem.} --- a proposition that in the following will find its application, when
we are led to see that philosophy must employ not merely a logically analyzing, but also a
psychological-historical method. Not all tasks are solved because ever more representations are
referred to the world of anthropomorphisms. When has one found the last anthropomorphism, the one that
men will not be able to free themselves from? It could be found only when, at some time, science's
first presuppositions and last results could lie before us. But there is here set an infinite task, on
which generation after generation must work, and which conditions the great horizon under which all
deeper-going research operates.

From the philosophical side it has long been inculcated that one does not get beyond the relation
between subject and object. Even if we could suppose ourselves to stand in the face of an absolute
objectivity, we should yet have to investigate what stage of subjectivity is hereby presupposed. Kant
and Fichte therefore rightly maintained that the strictly scientific apprehension of existence holds
only under the presupposition of a subjectivity that has the capacity for complete consistency and for
comprehending all possible experience. No actual human I can satisfy this demand. Every I (scientific
or unscientific) finds itself at a certain stage of development toward this ideal, which corresponds
to the ideal of truth or reality toward which all cognition is directed. It is not necessary to draw
forth this presupposition at every single question. Even the modern theory of relativity does not, as
a purely physical theory, draw forth this presupposition, but only demonstrates that every
determination of time, just as well as every determination of place, presupposes a definite relation.

In his characterization of the modern theory of relativity Émile Meyerson lays strong weight on its
objective essence. When nonetheless the presupposition of a possible observer (what Whitehead calls
the percipient event) appears, it is, according to Meyerson's conception, only in order to establish
that there is here talk of something that holds for every observer.\footnote{``\ldots une
représentation du phénomène valable pour tous les observateurs, quels qu'ils soient'' [a
representation of the phenomenon valid for all observers, whoever they may be]. \textit{Déduction
relativiste}, p.\ 78.} Such a pure I, as Kant and Fichte spoke of, should also precisely, in every
given situation, be capable of being represented by any I whatever, just as humanity at each given
time and at each given place is represented by single, definite men. The pure I is the human race
regarded as a great person who is constantly learning something.\footnote{Cf.\ Pascal
(\textit{Fragment d'un traité du vide}): ``Toute la suite des hommes, pendant le cours de tant de
siècles, doit être considérée comme un même homme qui subsiste toujours et apprend continuellement''
[The whole succession of men, through the course of so many centuries, must be regarded as one and the
same man who endures always and learns continually].} Science is the work of humanity, and Kant's and
Fichte's pure I is a personification of the human race, whose representatives the single researching
men are. When Einstein establishes that under certain conditions two observers cannot agree on a
determination of time, he himself takes a standpoint that lies nearer to ``the pure I'' than that of
any of the single observers, in that he sees the ground why these cannot agree.\footnote{On this see
\textit{Relation som Kategori} [Relation as a Category], pp.\ 74--76.}

Oswald Külpe\footnote{\textit{Die Realisirung}, Vol.\ III (1923), p.\ 427.} considers it inexpedient
to speak of a pure I, partly because such an expression can, as the philosophy of Romanticism shows,
easily be interpreted metaphysically, partly because the concept ``pure I'' plays no role in the
grounding of the fact that we stand in the face of a reality. --- To this it is to be said that the
danger of metaphysical interpretation ought not to deter one from setting up a concept that shows
itself to be necessary, and that the necessity in this case lies in the problem contained not merely
in the differences between contemporary observers, but also in the differences between the researchers
of the shifting times. He who sees the ground why different researchers (contemporary or from
different times) do not agree takes just as well a higher position, and so far designates a greater
approximation to the human race's ideal of cognition, as scientific cognition of reality designates a
higher position than the world-picture of the ordinary understanding.

The very subjectivity for which all possible human knowledge, both as regards compass and as regards
connection, could stand full and clear, would yet think under definite presuppositions, would be a
human subject ($S_m$). A pure $S$ (without the index that designates certain presuppositions) we
cannot think. But an idealized human subject ($S_m$) would have to possess an understanding of why
different empirical subjects ($S_\alpha\ S_\beta\ S_\gamma\ \ldots$), all belonging under $S_m$, would
have to come to different conceptions.\footnote{Cf.\ \textit{Den menneskelige Tanke}, p.\ 304.} $S_m$
is itself in constant development, and the full truth will consist not merely in an objective doctrine,
but also in an explanation of why not all researching spirits apprehend the matters in the way the
objective doctrine indicates.

In Leibniz's doctrine of monads\footnote{On this see \textit{Den nyere Filosofis Historie}, 3rd ed.,
II, pp.\ 135--149.} this standpoint is set forth in an ingenious manner, only that he introduced the
unnecessary assumption that each monad was a world closed in itself, in which the whole great world
mirrored itself.

\chapter{Organism}
\label{ch:2}

\hnum{4} When the relation between epistemology and the conception of life is to be investigated,
it will be instructive first to take the word ``life'' in the narrowest sense, that is, of organic
life. We do not, indeed, have in this the simplest example of the concept of the whole and of the
problems this concept sets. A globe, a mountain, a geological stratum, an atom are also wholes that
can provisionally appear as barriers to cognition, but can afterward show themselves to set more or
less significant tasks for it. The concept of the whole has precisely this double position. It
provisionally halts the logical analysis that is the predominant one in the formal sciences. In
particular the atom long stood as an absolute unity; only in the most recent time has it revealed
itself as a world of electrons, and significant investigations of what takes place in its interior
have thereby become possible and necessary. The ground that the concept of the whole (totality)
ought to be set up as a fundamental concept, a category, is precisely that it cannot without
further ado be led back to a relation of identity, but demands a peculiar series of investigations.
An intermediate link between identity and totality is formed by the concept of causality. For it
shows itself that what we call cause is never anything simple and single, but is a group of
positive and negative conditions; and it shows itself further that several different causal series
can meet, so that there appears a whole system of processes, of which the one conditions the other.
When such a system steps up as a given subject, it stands, epistemologically regarded, both as a
barrier (in the face of the application of purely formal concepts) and as a task (which demands
investigation of the conditions for its coming-to-be and subsistence).\footnote{Cf.\ on the ground
for setting up the concept Totality as a category: \textit{Den menneskelige Tanke}, pp.\ 218--222;
\textit{Totalitet som Kategori}, pp.\ 18--22; 48--52.}

There is an analogy between all wholes, from the simplest to the most complicated. Since the whole
called organism stands in many respects as typical of the role the concept of the whole plays in
our cognition, we hold first to it. By conception of life we thus provisionally mean scientific
conception of organic life. An investigation of this subject will be a good preparation for
discussing the conception of life in the sense of a view of life, that is, of the peculiar whole of
thoughts and moods that can form in men's minds under the influence of their experiences.

\hnum{5} From an epistemological standpoint it must be maintained that however strong an impression
of the whole an organism can make, it is yet unwarranted to suppose that it should beforehand be an
exception from the general natural (physical and chemical) laws. Every single organic process, and
within such a process every transition from link to link, must be attempted to be interpreted
according to the standpoints physics and chemistry take. This interpretation has, at least
hitherto, succeeded only to a certain degree, and especially the connection of the many different
processes into a whole-life that can maintain itself as such stands as a constant problem. This
problem is not solved by, as happens in the so-called vitalistic theories, attributing to the
organism as a whole a peculiar force (a ``vital force'') that should hold the elements together.
That would be to explain \textit{idem per idem} [the same by the same]. The concept of the whole
cannot scientifically be used to designate a cause or give an explanation. --- In the treatise
\textit{Totalitet som Kategori}, \S\,20, I have criticized so-called Vitalism. I here return to the
problem from a somewhat different standpoint.

It is characteristic of the manner in which the discussion of this question is conducted that the
contending parties mutually accuse each other of ``metaphysics.'' Those who maintain that in organic
life wholly new forces lay themselves before the day in relation to the physical-chemical domain
often find ``materialism'' in those who maintain that physical and chemical laws hold in organic
nature as in the rest of nature. In return, these latter often find ``spiritualism'' in the former.
Epistemologically the matter stands, in my conception, such that no boundary can be set once and for
all for the application of physical and chemical standpoints to the organic domain, but that on the
other hand such application cannot in fact be carried through. The problem of an organic whole's
relation to its conditions, to the single processes in whose cooperation the life of the whole
consists, can be solved neither by proclaiming Mechanism as the solution of all riddles nor by
letting the whole ``itself'' intervene as a \textit{deus ex machina}. It is by this posing of the
problem that the investigation of the concept of organism will show itself to be an analogy to the
investigation of the concept of personality and of views of life and of the world.

\hnum{6} As researchers who with more or less consistency try to hold this line, one must name J.\
S.\ Haldane, whose book \textit{Mechanism, Life and Personality} (1913) I was already able to
discuss in the treatise on Totality as a Category, and from the most recent time E.\ S.\ Russell
(\textit{The Study of Living Things. Prolegomena to a Functional Biology}, 1924). This last author
seeks in an interesting way to distinguish between biology and physiology. The method of physiology
is, he maintains, purely analytic; it resolves the complex phenomena of life into elements that
stand in physical and chemical reciprocal action with one another. Biology, on the other hand,
investigates life in its immediate appearance. This appearance, with the unity and coordination it
displays, must be taken as a fact, however much its elements or components can be analyzed in
detail. Biology studies life in its free unfolding and finds precisely the urge to free unfolding
to be the characteristic of life in opposition to the lifeless. This urge Russell will designate
with the Greek word \textit{hormé} (conatus).\footnote{Cf.\ on Spinoza's concept of conatus:
\textit{Spinoza's Ethica}, pp.\ 56; 61.} A living being should first and foremost be studied in its
natural surroundings, to which its movements are more or less adapted. This is the task of
functional biology. Already in the phenomena of growth, that is, in the organism's successive
coming-to-be, the necessity shows itself of taking account of the organism as a whole. How does it
come about that the bone-cells are formed at the right time and in the right place? And why is the
formation of bone necessary at all? Questions such as these can be answered only when the single
process is regarded in its relation to the organism's life as a whole, especially to its life in
the future. There shows itself here a historical solidarity between the single process and other
processes within the same whole. There is thus here need of synthesis, not merely of analysis, when
understanding is to be won.

Russell compares physiology's relation to the experiences open-air biology can make out in nature
with the relation of aesthetic criticism to the work of art. Such a comparison already underlies
Kant's \textit{Kritik der Urteilskraft}, which treats both the concepts of art and genius and the
concept of organism. It is the involuntary and complex phenomena in nature and in the world of
spirit that are here set side by side. They stand as unintelligible, perhaps paradoxical, and it
perhaps never succeeds for the analyzing research to find an exhaustive explanation of them.
Open-air biology and experimental physiology often stand here in sharp opposition to each other. Out
in life, as it stirs in gardens, marshes, and forests, the forces can unfold themselves in many
different directions; but the analyzing physiology is inclined to acknowledge only what can be
demonstrated experimentally, that is, under certain definite, delimited conditions. And yet ---
what Russell does not sufficiently bring out --- experiment and analysis are, after all, the
purgatory every truth must go through. The matter is that in free nature it will as a rule not be
possible to demonstrate the simple and immediate relations of succession, since these are not
accessible to direct apprehension. Kurt Lewin (\textit{Der Begriff der Genese}, 1922) has, I think
rightly, maintained that it is for the most part through inference from likenesses between different
phenomena that one reaches the assumption that they are successive links in a developmental series.
He thereby explains why in freely living animals one can so seldom infer lawfulness in inheritance.
The demonstration of immediate series of succession stands precisely as the verification of the
suppositions open-air biology forms about the connection between forms of life. It then depends on
whether one can form such series in which every link stands as the immediate continuation of
another, as in the series of stages of growth in the same organism and in the succession of
generations in the same group (in so far as this succession can constantly be followed). By the word
``gen-identity'' (genetic identity) Kurt Lewin designates the identity between subjects that rests
on their immediately succeeding and continuing one another, belonging to the same genetic
series.\footnote{Lewin's series of gen-identity correspond most nearly to what I have called
identically varying series of difference (\textit{Den menneskelige Tanke}, p.\ 167).}

If Kurt Lewin is right, the necessity shows itself of experimental physiology for the revision of
the results open-air biology (functional biology) supposes itself to have reached. The two
directions of research must constantly supplement each other. Whole-apprehension and analysis do
not exclude each other; the former constantly furnishes material for the latter, and the latter
will scarcely ever be able to resolve organic life into a series of purely mechanical links. An
analogy that here lies near (besides the one used by Russell) is the relation between geometry and
arithmetic. When we draw a line, we constantly add piece to piece, but we cannot establish that one
addition is just as great as another. Nonetheless the series of geometrical additions and the
series of numbers are analogous, such that to every transition within the one series there
corresponds a transition within the other. For the understanding of organic life it would be the
ideal if to every new mode of view in functional biology there could be demonstrated a mode of view
agreeing with it in experimental physiology.

\hnum{7} The standpoint of the whole thus does not stand in complete opposition to the mechanical
conception. They are related as synthesis to analysis. And the one cannot exclude the other. But the
standpoint of the whole cannot in and for itself be used for explanation. Not merely in biology,
where vitalism and teleology have sometimes come forth as definitive solutions of riddles, also
elsewhere has one wished to employ the standpoint of the whole against the mechanical conception of
nature that rests on strict analysis. After his grand demonstration of the connection within our
solar system --- of the manner in which the various globes' masses, distances, velocities, and
densities stand related to one another, and in which the planets move in circles around the sun
instead of, in consequence of the law of gravity, falling down into it --- Newton declared that
this whole wonderful connection (\textit{elegantissima compages}) could not have arisen according to
mechanical laws.\footnote{\textit{Den nyere Filosofis Historie}, 3rd ed., II, p.\ 208.} But analysis,
and with it the attempts to explain the coming-to-be of the joining-together, came later, even if
there constantly remain unsolved riddles.

The opposition between totality and mechanism, synthesis and analysis, thus comes forth again and
again for our cognition. In the following sections we will meet them in the domain of spiritual
life.

\hnum{8} E.\ S.\ Russell does not let himself be content with the consideration hitherto set forth.
For him the significance the standpoint of the whole has becomes clear only when one employs, as an
analogy (and something more than mere analogy!), a glance at the whole that human soul-life presents,
``so that at the foundation of living beings' material appearance there lies an activity of a kind
similar to (\textit{similar to}) that which we know as the innermost reality of our own life.'' On
this point he diverges from Haldane, whom he otherwise regards as one of his predecessors. Our own
experience is, he supposed, the only measure of life-activity we have, and whether we will or not,
we must apprehend all other life-activity after this model. And it is then especially the tendency
or striving toward a goal appearing in our soul-life that cannot be dispensed with if functional
biology is to be able to solve its task. --- In a lecture which Russell, the year before his book
appeared, gave on what he calls ``Psychobiology,''\footnote{\textit{Proceedings of the Aristotelian
Society}, 1923.} he expressed himself more positively about his method and even calls it outright
psychological; he also states here that the problem is probably insoluble, but that we are referred
to going at it from two opposite ends of the scale, the material and the psychological.

It is indeed so, and it has already been shown by Kant, that we dispose of only two analogies for
interpreting what organic life is: the analogy with mechanical causality and the analogy with
psychic striving after a goal. The choice stands in the end between mechanism and personality. But
just as unwarranted as it would be to declare the organism a pure mechanism, just as unwarranted is
a dogmatic application of psychological concepts beyond the domain where experience directly compels
it. Only with great caution and with all possible reservation can one speak of unconscious
soul-life. Even in the elementary expressions of soul there must be hesitation about applying
psychological representations that are throughout drawn from the more developed life of
consciousness.\footnote{Cf.\ \textit{Den menneskelige Tanke}, pp.\ 1--19.} And since now, according
to Russell himself, the concept life is more elementary than the concept soul, nothing in the end is
gained by setting up the concept ``psychobiology.'' It is better to hold to the expression
``functional biology.'' This must set out to show what a \textit{compages elegantissima} the organism
is, and then leave it to physiology to attempt a mechanical explanation of it. The next question will
then be whether such a mechanical explanation can exhaust all that functional biology can lay before
us. One must in science not merely go from a synthesis at hand to analysis, but must be able to go
from this analysis to a new synthesis that makes the first synthesis at hand intelligible.

When Russell makes use of psychological analogy, it is not because he supposes that soul-life appears
as a wholly new and foreign principle in relation to the mechanical connection of nature. He is not a
dualist, but attaches himself most nearly to the Spinozistic-Leibnizian conception, according to
which the psychic is potentially present even in the most elementary natural processes, although it
first actually emerges at a certain stage within the development of organic life. It is now also
clear that, if one assumes a continuity in existence at all, one must assume preceding possibilities
for the new forms of existence that arise. But here too it is analogy that guides us, and which may
not be applied dogmatically. The natural philosophy of the seventeenth century apprehended rest as
equilibrium between acting forces, in order thereby to get beyond a sudden leap at the transition from
rest to motion or conversely. But that rest is thus only potential motion can be established only by
the fact that motion sets in as soon as the equilibrium between the forces is cancelled. A
possibility that never becomes reality is an impossibility. The concept of potentiality therefore
constantly contains a problem, not a solution. Functional biology has in its domain the significance
of setting problems that would otherwise not be set. Something new does in fact arise in the world,
without its being immediately seen to be a transposition into new form of something that previously
was. Thus life emerges in relation to the inorganic (if it is not, as some suppose, just as original
an element in nature as the inorganic). And thus soul-life arises in relation to organic life, and
within soul-life the more complex forms of soul-life arise as something new in relation to the more
elementary forms. One can here often demonstrate only a chronological relation, not a relation of
production. G.\ H.\ Lewes rightly spoke of emergents as different from resultants, and in his address
at the British Association's meeting in 1921 Lloyd Morgan took up this concept, speaking of ``emergent
evolution.''\footnote{\textit{Advancement of Science}, 1921.} He thereby wished to express that the
new must be taken as such, not without further ado dogmatically seen as inexplicable. Unexplained it
is provisionally, and perhaps forever. It stands as a whole that waits to be analyzed. William James
had already employed this concept, when for him the expression ``free will'' meant only something new
that emerged in the domain of the life of will, without regard to whether an explanation of it was
possible or not.

\hnum{9} Without wishing to go into the great question of the origin of life on earth, the American
chemist Henderson\footnote{\textit{Bulletin de la société française de philosophie}, 21 January
1921.} has maintained that the origin of life was prepared by the fact that our planet's development
very early entailed the formation of an atmosphere containing water and carbonic acid besides other
substances, and thereby produced the most important conditions for organic life. Life has thus not
entered into a world that was foreign to it or offered only hindrances. And in this Henderson finds a
warrant for regarding existence from a teleological standpoint, in that the more elementary forms of
existence, according to their own laws, form a favorable foundation for higher forms of existence.
The expression ``teleological'' Henderson is not himself satisfied with, just as little as with the
expression ``\textit{finalité},'' which he used when he set forth his ideas in the French
philosophical society. It was here brought out during the discussion, from the chemical side, that the
conditions on the earth's surface could in and for themselves have made possible an origin of wholly
other forms of existence than the organic, so that it was not warranted to speak of preparation or
harmony between elementary and higher forms of existence. From the philosophical side it was
inculcated by Meyerson that the expression teleology or finality ought only to be used when there can
be talk of a conscious ordering of conditions to serve a goal, and that one must not appeal to
non-empirical causes so long as one cannot furnish proof that every possibility of finding natural
causes is excluded.

Despite these warranted objections against the train of thought set forth by Henderson, the fact yet
remains that life, however it has otherwise arisen, has found favorable conditions present. Otherwise
it would naturally have vanished again. This acquires significance for the illumination of the
subjects with which the present treatise occupies itself. However great the problems that the
formation of wholes at all stages and under all forms sets our cognition, there is yet every reason
to continue the investigations of the relation between the wholes, their elements, and their
conditions. There is perhaps an irrational relation present here, as between the diagonal of the
square and its side; but there can, as here, be found more and more decimals. And for our whole
conception of life and the world the subsistence and unfolding of the whole acquires a decisive
significance. For to the concept of the whole the concept of value is (as I have earlier shown)\footnote{\textit{Den
menneskelige Tanke}, pp.\ 238--240; \textit{Totalitet som Kategori}, Ch.\ 5.} closely bound, in that
it is everywhere conditioned by the relation between the subsistence of wholes and their elements and
conditions.

\hnum{10} By this consideration a light can be cast on an essential point in Kant's philosophy. As
indicated in the foregoing, for him the opposition between mechanism and organism was not fundamental,
but was due to our cognition, which is here referred to employing two different methods --- partly to
seek to go from the parts to the whole, partly to seek to go from the whole to the parts --- without
being quite able to traverse either of these two ways. It was Kant's conviction that it is the same
order of things that underlies mechanical connection and the coming-to-be of organic wholes. The
beauty and purposiveness of nature we therefore have no right to regard as accidents, but we must in
them, just as well as in the causal connection that appears in the phenomena, see expressions of the
nature of existence. Such a conception already underlies Kant when, in his youth, he disputed Newton's
right to declare a mechanical origin of the solar system's harmony impossible, and when he later, in
his principal work, rejected the teleological proof of God's existence. And in his \textit{Kritik der
Urteilskraft} (1790) he teaches that nature comes to meet the human spirit's need and apprehension by
making beauty and purposiveness possible.

By this train of thought there is really cancelled the sharp opposition that Kant himself, in
\textit{Kritik der reinen Vernunft} (1781), had presupposed between human cognition and ``the things
in themselves.'' In the beautiful and purposive phenomena of nature ``the things'' come to meet us;
why then should they not come to meet us when we strive after cognition? --- Kant had, by his
opposition to Dogmatism, been led to a sharp distinction between what lies in the nature of human
thinking and what lies before us as given and must be received as it is. It is, after all, the
self-activity of thinking that makes science possible. But Kant set out from the assumption that human
thought was like a stranger in the world and could not expect support or confirmation from existence
itself, which he called ``\textit{Ding an sich}.'' Yet he assumes in his principal work that existence
itself is the cause not merely of the ``matter'' a cognition has to work up, but also of the ``forms''
in which the working-up is undertaken. But if it is so, then existence comes to meet us here too! ---
The train of thought on which \textit{Kritik der Urteilskraft} is built is thus, when one looks
closely, not essentially different from the one that underlies Kant's epistemology, and truth has,
like beauty and purposiveness, its ground in the nature of existence itself.\footnote{This conception
of Kant's philosophy I have already set forth in my treatise \textit{Kontinuiteten i Kant's
filosofiske Udviklingsgang} [The Continuity in Kant's Philosophical Development] (Vidensk.\ Selsk.\
Skr., 1893). It is briefly presented in \textit{Den nyere Filosofis Historie}, 3rd ed., III, pp.\ 59;
106.}

\chapter{Personality}
\label{ch:3}

\hnum{11} With the concept of personality, as with the concept of organism, the thought of the
whole steps into the foreground, and with it there also appear the problems this thought sets in
motion. In the history of psychology there have appeared conceptions according to which human
spiritual and soul-life was explained as having arisen by a summing-up of the single sensations,
representations, feelings, and drives that one could demonstrate. One supposed that development in
the psychic domain went from the single and simple to the complex. Such a conception is already
refuted by the fact that it is first at a later stage of development that one can definitely
distinguish between different psychic elements. In the child and the savage there is already
concentration, in the form of reflex and instinct. Here one cannot distinguish between different
sides of conscious life; they work immediately together. And if any of the elements that
psychological analysis supposes itself able to distinguish is to be regarded as original, it would
have to be willing, this word taken in its widest sense. As one of the most primitive psychic
phenomena one has, especially from the psychopathological side, pointed to ``motor
whole-phenomena'' (Gadelius), where an influence immediately leads to movement. Reflex and instinct
are peculiar, more or less complex forms of such whole-states. The urge to movement is from the
first ruling and forms the constant but dim foundation, out of which, and in whose service, all
cognitive life and all feeling-life can later develop. And in reciprocal action with these the
original urge to movement, which is at the same time an urge to life, an urge to subsistence and
unfolding, can develop into more special and more peculiar forms. When an urge is combined with a
representation of a goal, it can be called a drive. A collision between different urges or drives can
lead to an inner struggle that can be clarified through what we call deliberation, during which now
one, now another drive has the preponderance, until a single drive, with more or fewer
after-effects of conflicts, conquers --- in looser form as intention, in firmer form as
decision.\footnote{Cf.\ my treatise on \textit{Begrebet Villie} [The Concept of Will] in
\textit{Mindre Arbejder} III.}

In this whole course of development, which forms a series whose links cannot be interchanged, the
life of representation and feeling is essentially serving, however significant its cooperation may
be. --- What observation of soul-life during its development from child to adult, or from the
``wild'' life to the life of culture, thus shows is confirmed by what animal psychology and
sociology teach us. Thorndike has through his experiments with animals shown that development in
them goes from undifferentiated states to discernment and generalization, and sociologists such as
Tönnies and Lévy-Bruhl have shown that man's soul-life at earlier stages has a complex character,
without our usual psychological distinctions being applicable there.

When now not merely urge or drive stirs in different directions, but when also cognitive life and
feeling-life can stir with a certain independence of the original urge to life, then the question
becomes whether soul-life can continue to be a whole. The character of the whole rests upon the
synthetic energy with which the elements of soul-life are worked together. A chaos of the elements
designates the transition to mental illness. And one can therefore distinguish between synthesis as
the fundamental form of all soul-life, and the content that is comprehended with greater or lesser
energy and with greater or lesser success. Already in my \textit{Psychology} (1882) I have therefore
distinguished between a formal and a real self. The former announces itself in the synthetic energy,
the latter in the qualitative peculiarity of the content that is comprehended. When I so strongly
brought out the concept of synthesis, it was because it was then a matter of the unity and coherence
of conscious life in the face of the psychological atomism that prevailed in the English
psychology of association, which in other respects has its great merits. I pointed already then to
the fact that in a genetic presentation one would have to begin with states that were most nearly to
be designated as states of will, when one took the word willing in its widest sense. Unfortunately I
did not come to rework my Psychology from this standpoint.

Peculiar problems arise from the fact that the soul-content cannot always be comprehended in a
harmonious manner. As already mentioned, there can take place conflicts between different urges and
between conflicting drives, so that it is a matter of working them together into a whole, which can
perhaps happen only by the repression of single refractory elements. Spiritual development begins,
like bodily development in the embryonic state, often from scattered points. On this sporadic
character Sibbern laid great weight in his Psychology. It is this that sets us our tasks. ---
Purely theoretically there arises the task of finding a bond of unity between our different
observations, and in analogy with this there is set us the task of building up a character out of
the chaos of drives and impulsive urges that more or less stir in each of us. In the so-called
hypnoid individuals a disposition makes itself felt to a cleaving of soul-life, in that there stir
tendencies to highly different ways of taking life. A single such tendency perhaps gains the upper
hand, and there can then take place a whole-formation around it as center-point. It can even, as
especially Flournoy\footnote{Cf.\ Claparède, \textit{Théodore Flournoy. Sa vie et son œuvre} (1921),
pp.\ 56--65.} has shown, be salutary that such a partial totality develops, since it can thereby
become possible that a difficult task is solved or a dangerous impulse overcome, although a complete
gathering of all the mind's powers is not reached. Alienists often cultivate in their patients a
single tendency that has particular value, in order through this to repress other, less fortunate
partial totalities. Indeed, one can go a step further. No real development of character is possible
without some tendency or interest coming to rule, whatever value it may otherwise have. It is
incorrect when we speak of our ``having'' certain qualities. Our essential qualities we do not
``have,'' but we ``are'' them. Napoleon is said to have said that he did not ``have'' ambition, but
that ambition penetrated his whole being, just as the blood streamed through his whole body.

In Gustaf Fröding the inner struggle between great spiritual oppositions for a time led to madness.
His letters, especially those to his sister, give about this information of great psychological
interest. He speaks of the terrible split between delicacy of feeling and coarseness, between shyness
and defiance, that prevails in him. And to this comes the uncanny self-criticism that almost makes
his artistic creation an impossibility. In his despair he then sometimes puts on a buffoon's costume
and performs the corresponding role so well that he almost goes mad. If to this came an absolute
dryness and emptiness, where it was possible to feel neither love nor hate, neither joy nor sorrow,
it was no wonder if madness seized him. And yet under all this his now exalted, now sensuous fancies
could form themselves into coherent wholes. Sometimes he had a strange sensation of being another
person, now the one, now the other, and that to such a degree that not merely his interior, but also
his facial features and his figure corresponded to the person he supposed himself to be. He rightly
adds that what he has thus experienced must be akin to what the so-called mediums experience. And
despite everything he had in brighter times a conviction that these experiences did not merely mean
illness, but pointed to ``a process of development of the personality, long prepared by preceding
broodings.''\footnote{Gustav Fröding, \textit{Bref} (1915), pp.\ 206--210; 219\,f.} Full personal
unity the great poet did not, indeed, reach. But in return he has depicted the opposition between
light and darkness, and between dream and reality, as beautifully and penetratingly as no other poet
of our time. It is easy enough to gather one's life-content in harmony when it presents no sharp
oppositions. The worth of harmony rests on how great disharmonies it has had to struggle with.

Not merely simultaneous oppositions can cleave the real I and make synthesis difficult. Successively
appearing oppositions too can have such an effect, especially while the transition is going on. Plato
halted at this difficulty on account of the sharp distinction he made, in the older form of his
doctrine of ideas, between the different ideas --- that is, between the different predicates that can
successively belong to the subject, or, let us say, the different content with which one and the same
I can gradually be occupied. Strangely enough, he did not, with this problem (in the dialogue
Parmenides), think of the fact that he himself, in his ingenious doctrine of Eros (in the Symposium),
has depicted a psychological course of development in the course of which the soul gradually acquires
new content.\footnote{\textit{Bemærkninger om den platoniske Dialog Parmenides} [Remarks on the
Platonic Dialogue Parmenides] (Vidensk.\ Selsk.\ filos.\ Meddelelser, 1920), pp.\ 44--50.} He lays
less weight on the psychological difficulties, the pain that can arise at the transition from link
to link in this development, than we would do,\footnote{This remark I owe to my wife, who in
conversations has often criticized Plato's psychological doctrine of evolution in the ``Symposium''
from this standpoint.} but has as a psychologist had a faith that such a development was possible. It
is first in more recent time that the psychology of transitions, of crises, has become an important
subject for research, and we can often need Plato's faith in order not to doubt that bridges can be
built over psychic abysses.

A complete comprehending of the life-content is not possible for any personality. There constantly
emerges new content, and with it new tasks for comprehending. An absolutely gathered personality is
an ideal, which yet does not exclude that, for the attentive eye, at every stage and in every form of
psychic life, there are discovered intimations of, and approximations to, what such an ideal demands.
This holds both in the involuntary and in the fully conscious soul-life, although both we ourselves
and others best understand what makes itself felt in the ruling feeling and in the constant thought.
What at the moment, or hitherto in most moments, has occupied the center-point of consciousness need
not be what is in reality most deeply grounded in our nature. Much experience and reflection is
required in order to become clear about our own real I, and it becomes really possible, indeed, only
when opposition to other I's is noticed. As regards the formal I, which consists in the synthetic
energy that holds the soul-content together, the concept of it naturally becomes possible only in the
face of a comparatively advanced intellectual development. From the first one cannot distinguish
between form and content. The word I is naturally not decisive. For children it is provisionally a
name, and it can happen that they come into strife about which of them has the right to bear this
name. Philosophers have often been accused of laying great weight on the use of the very word I. But
even Fichte, who speaks so much of ``das Ich,'' and who has often (so recently by an excellent Danish
philologist) been mocked for it, is clear that the thought ``I,'' by which he especially means the
formal I, can first arise when a certain intellectual development has taken place. ``So long as I,''
he says,\footnote{\textit{Rückerinnerungen}, \S\,36.} ``live wholly practically, am wholly life and
acting, I know only my feeling, urge, willing, and so forth in their changing, but I do not expressly
know myself as the unity and principle of these different determinations. Only when I raise myself
above the reality of these different acts and, with abstraction from their difference, merely
comprehend them as all belonging in me, does the consciousness of the unity as the principle of those
manifold determinations at all arise; and this product of our abstracting and comprehending thinking
is what we call our soul or spirit.''

\hnum{12} The distinction between the formal and the real in personality is of significance for the
understanding of the relation between epistemology and psychology. Epistemology singles out the
formal and looks away from the real content that varies from individual to individual and, in the
same individual, in different periods. The concept of synthesis is a fundamental concept for both
disciplines. All scientific concepts bear its stamp. It is no wonder, since scientific work is
exercised by personalities. But it is at the same time a work that leads beyond everything
individually personal, in that it sets out to win an understanding that can become common to all
whose capacity of thought reaches a certain development (cf.\ above, \S\,3). An understanding would
be impossible if all subjects at hand formed a chaotic series of difference --- that is, a series
whose links were all equally different among themselves, could be placed in any order whatever in
relation to one another, and each be thought away. It will indeed never be possible to establish that
such a series really lies before us. But it can be used to illuminate the nature of the work of
thought. Already the ordinary human understanding seeks beyond the chaotic, where it to a certain
degree announces itself, and involuntarily connects the subjects in one way or another. Science
continues this work through strict testing of the single subjects and of the ways in which they can
be connected. Through a series of stages, of which the most important are more closely characterized
in \textit{Den menneskelige Tanke} (pp.\ 160--173), science seeks to reduce all differences to being
hidden identities. Pure identity stands so far as science's ideal. But the very concept of identity
is, like all concepts, the fruit of a synthesis, in that it designates a relation between subjects
that hitherto stood as different. Where such synthesis succeeds, science has won a victory.

\hnum{13} There have been times when one reveled in identity as the summit of all science. To lead
all subjects back to one subject, which stood as identical with itself and therefore needed no
explanation, while it at the same time underlay all other subjects, stood as the great ideal, which
one sometimes did not doubt one had reached. A boundary for the carrying-through of this ideal appears
already in the qualitative differences, with their varying in degree, time, number, and place. It is
here the rational analogies (see above, \S\,2) that step into the place of identity. But a
particularly significant barrier is here designated by the concept of the whole, which finds
application where different qualities in different relations of degree, time, number, and place appear
in experience as connected in more or less firm coherence. Such a whole is the organism, and such a
whole we have also in the personality --- the personality in which science itself has its origin. The
remarkable thing is thus that by personal work there have been produced thoughts that seem to threaten
the very reality of personality. For when identity is the expression of truth, personality cannot be a
true concept, partly because there are many personalities, partly because every single personality is
a peculiar factual synthesis of different subjects and elements. Identity and totality stand here
apparently as hostile powers over against each other. When totality is nonetheless taken up as a link
in the series of categories, it is first and foremost because this concept draws attention to peculiar
tasks for thinking.\footnote{See further \textit{Den menneskelige Tanke}, pp.\ 218--227, and the
treatise \textit{Totalitet som Kategori}.} As regards personality, the task becomes first that of
apprehending the single personality in its peculiarity, living oneself into it, provisionally seeing
things from its standpoint --- in other words, demonstrating the particular form of existence, different
from other forms, that here steps forth. It is further a matter of investigating how all elements in
the personal whole are in their place there, so that they each contribute to the whole's subsistence
and development. Even disharmonies can be in their place, if without them a fullness of harmony as
great as possible for the whole cannot be reached. Already here analysis is necessary. But as a third
task there appears thereafter that of investigating the conditions for such a whole's coming-to-be and
continued subsistence. That happens by the psychological-historical way, and this investigation acquires
a genetic character, in that it sets out to find the intermediate links through which the personality
has reached the characteristic form in which it appears. This task is at bottom never fully soluble. The
personal stands as an irrational magnitude that can be determined only approximately by particular
numbers. It is in relation to the pure principle of identity that personality stands as something
irrational (in the sense stated). But the concept of totality, and with it especially the concept of
personality, gives occasion to another concept of irrationality. From this concept the irrational
becomes that which does not satisfy the fundamental condition for a personality's subsistence in its
full peculiarity, and which therefore can with difficulty, or not at all, be thought as a lasting
element in it. Especially in the investigation of a personality's genesis many riddles can remain
unsolved. But just as little as when the talk is of organic life does the whole itself contain an
explanation. When one appeals to one's personality, one is at the end of one's cognition, and the same
holds when one would explain another man's personality by --- his personality. There is a psychological
vitalism just as there is a physiological vitalism. There can arise a collision between two different
interests here: that of reaching a full characterization of a personality in its individual peculiarity,
and that of reaching an understanding of all the elements in this personality and of the manner in which
they have joined together into this whole. The two interests do not at bottom conflict with each other,
but it is seldom that they make themselves equally strongly felt in the same researcher.

\hnum{14} In his interesting books on Irrationalism and Individuality\footnote{\textit{Der
Irrationalismus} (1922). --- \textit{Philosophie der Individualität} (1920).} Müller-Freienfels uses
throughout the word irrational of what does not go in under the thinking that works with the principle
of identity. From such a standpoint the concept of totality becomes irrational. So especially all
individuality. Science (historical research and psychology) then gets the difficult task of making what
is in itself irrational an object of rational understanding. --- In opposition to this conception I
think that one must apprehend totality as an independent category which, without denying the principle of
identity that no thinking can dispense with, grounds a peculiar measure, in that it summons one to
investigate the relation between a whole at hand and its elements. With this it follows that as irrational
there now comes to stand everything that cannot become a link of, or a condition for, a certain given
whole, and yet acquires a promoting or inhibiting influence on it. --- It is yet seen that the concept of
totality also in Müller-Freienfels lies behind certain forms of irrationality. So when it is declared
(\textit{Der Irrationalismus}, p.\ 218) that ``by the concept irrational'' there is to be designated ``a
positive, a factual life-power,'' such as can be traced in art, morality, and religion. If one once
acknowledges such a life-power, then there is thereby introduced, I think, a new concept of
irrationality, which is determined not merely by the opposition to the principle of identity, but also by
the opposition to what the principle of the whole demands in a given connection. One can place oneself
purely and simply on the standpoint of the principle of identity and declare everything that cannot be
fully grounded from it to be ``opinions of mortals'' or ``anthropomorphisms.'' But in the chaos that thus
yet remains there appear groups, wholes, which demand reflection that cannot content itself with the
principle of identity.

The concept of the whole is the necessary presupposition for the concept of value. Of value one can
speak only from the standpoint of a whole. Everything acquires value that serves a certain whole's
subsistence and unfolding. Since Müller-Freienfels does not have the concept of the whole as an
intermediate link between the concepts identity and value, it is understood that he denies that there can
be a rational ordering of values. ``\textit{Die Werte bauen sich nicht auf Gesetze auf}'' [Values do not
build themselves up upon laws] (\textit{Philosophie der Individualität}, p.\ 204). This is in a certain
sense correct. One can speak of value only in relation to a certain given whole, but the latter's
conditions of life will be able to be arranged in a series whose law depends on the character of the
whole's need and circumstances.\footnote{\textit{Den menneskelige Tanke}, pp.\ 346--349. ---
\textit{Totalitet som Kategori}, Ch.\ 5.} That great problems can here come forth is a matter for itself,
which we will come to mention in the next chapter. A new group of irrationalities arises, after all, only
when the concept of totality is laid as a foundation beside the concept of identity.

\hnum{15} Just as work is done to explain the organism as a mechanism whose single links can be
understood by the physical and chemical way, so it is often attempted to explain personality as a product
of more elementary forces. It is in both cases the concept of the whole that is overstepped.

One has supposed that what we call personality, and indeed all psychic phenomena, are mere effects, or
perhaps side-effects, of the physiological processes. About attempts in this direction I have often had
occasion to express myself, most recently in \textit{Begrebet Analogi} (112--119). According to the
conception I have sought to ground, the relation between physiology and psychology is an example of
rational analogy, whereas neither of these disciplines can be grounded by pure deduction from the other.
Thereby both materialism and the ``psychobiology'' mentioned in the foregoing fall away. By means of the
rational analogy between the two disciplines, physiological knowledge and analysis can in many ways partly
strengthen, partly correct what psychological investigation has led one to assume. One can thus (as
especially Bekhterev has shown) distinguish in the nervous system and its functions a central and a
peripheral. Central organs can be demonstrated in which after-effects of earlier influences can constantly
be, as it were, preserved in a kind of ``collecting-basin'' for the experiences most important for the
organism, so that momentary influences do not become wholly decisive for the organism's reactions. But
when Bekhterev calls these central organs ``the personal sphere,'' it is clear that he builds on an analogy
drawn from psychology. Through psychological observation a difference between a central and a peripheral can
very well be demonstrated, without the use of any physiological insight. The central, which I usually call
the real I, we regard as more characteristic of the personality than the sensations and impulses emerging
in single moments. --- How far the rational analogy between physiology and psychology can be carried
through, only continued experience can decide. For the present there is no ground to believe other than
that all psychic states and functions are bound to physiological states and functions, even if one cannot
always demonstrate these. Beforehand no boundary can be set for what we can hereby learn about the
conditions of our soul-life.

Another attempt to establish the secondary character of the concept of personality was made within
psychology itself, in that the older English school saw, in the mechanics of the association of
representations, the way along which the personality as a whole came to be through reciprocal action
between the single psychological elements. But it shows itself that all forms of the association of
representations can be led back to a tendency to gather the content of representation into a whole, and this
whole-association bears witness that the urge and capacity for comprehending constantly expresses itself.
It is the task of psychology to demonstrate a synthetic energy as active in all the phenomena of soul-life
from the most elementary to the most complex. What we call personality is only a marked form of, and result
of, such psychic energy.\footnote{Besides referring to my \textit{Psykologi} [Psychology], I can here refer
to the investigation of the origin of ``total feelings'' that I have made in my book \textit{Den store
Humor}.}

Neither as regards the personality nor as regards the organism can the synthesis of life be completely
explained by the way of analysis. There will here constantly stand a struggle between two different
principles: the one, which sets out from the elements and would explain the whole as mere composition of
them; the other, which seeks to find the stamp of the whole everywhere, even in the simplest organic and
psychic phenomena.

\chapter{The View of Life}
\label{ch:4}

\hnum{16} The science of spirit becomes possible only by the fact that personal life unfolds itself
freely and naturally. Only thereby does the science of spirit get subjects to discuss. This holds
especially for the philosophical disciplines (epistemology, psychology, ethics, and the philosophy of
religion). There must come attempts at science in order for epistemology to get a subject to
investigate. Soul-life must unfold itself involuntarily in order for psychological observation and
analysis to become possible. And an ethical and religious life must develop in order for ethical and
religious problems to be raised and discussed. And such involuntary unfoldings will, as experience
shows, lead to the formation of spiritual wholes. As already mentioned, a law of the whole underlies
all association of representations. Most acts of recognition consist, especially when it is a matter of
rapid reflection, in a part (e.g.\ a single quality) being immediately recognized and the rest of the
subject's qualities being involuntarily presupposed. By such partial recognition sense-illusion easily
becomes possible. In such a case an involuntary synthesis is undertaken, which can be corrected by a
subsequent analysis. It is a well-known psychological phenomenon that can serve as a simple example of
the relation between involuntary synthesis and a following analysis. As Lévy-Bruhl has shown, there is
a primitive tendency to make part into whole --- that is, to regard the presence of a single element as
a symptom that a complex of different elements is present. A characteristic example of what Lévy-Bruhl
calls participation is that savages believe that by having an image of a thing one has the whole thing
in one's power.\footnote{In \textit{Begrebet Analogi}, pp.\ 7--23, I have discussed this more closely
and also sought to show that here it is a psychological phenomenon that can give an explanation of a
sociological phenomenon, not the reverse.}

In phenomena such as those now indicated we have a starting-point for the understanding of the
coming-to-be of views of life --- perhaps, as will later be mentioned, also for the understanding of
metaphysics and religion. For the present we hold to the individual ways of apprehending and taking
life, that is, to syntheses that arise involuntarily in single individuals in a certain independence of
science and religion. It was a peculiarity of the nineteenth century that it was maintained by a number
of eminent thinkers that in all faith and view of life it depends first and foremost on genuineness ---
that is, on the faith's having really proceeded from personal need under the influence of the
life-experiences that had fallen to the individual's lot. Carlyle opens the series with his \textit{Sartor
Resartus} (1836); here in the North there appears shortly after (but without the influence of Carlyle)
Erik Gustaf Geijer, and then Danish thinkers such as Poul Møller, Sibbern, and Kierkegaard.\footnote{More
on this in \textit{Ledende Tanker i det nittende Aarhundrede}, pp.\ 137--142.} In Germany Nietzsche's
lasting significance has shown itself to rest on his demand for genuineness and independence in the view
of life. Dilthey too has, on the basis of his thorough studies of philosophical border-questions, worked
in the same direction, and in more recent time Georg Simmel is especially to be named, who in the
masterly little work \textit{Der Konflikt der modernen Kultur} (1921) has with great clarity illuminated
the unrest of the time. While in France Bergson, from his spirited demonstration of life's urge to
unfolding (\textit{élan de vie}), at once leaps over into metaphysics, Simmel remained standing on the
threshold. Even if it should be his early death that hindered him from attempting a journey on the other
side of the boundary, it is yet characteristic of him that he hesitated so long. It was his studies of
Schopenhauer and Nietzsche that led him to lay so great weight on the view of life's springing from the
innermost source of personal life. In his above-mentioned work he shows in an interesting way how so
many of the forms in which life had hitherto arranged itself are in dissolution, and how life, in its
impatience that no new forms offer themselves, stirs in many men with hostility toward all form. It is
indeed characteristic of spiritual life in Germany at present that one yearns after a view of life that
first and foremost bears the personality's own stamp, so that science and religion, if they work along,
stand as subordinate elements. In one's investigation of the different views of life one therefore lays
particular weight on the psychological foundation that ultimately bears it. Karl Jaspers has treated this
subject in detail (\textit{Psychologie der Weltanschauungen}, 2nd ed., 1922). He declares Kierkegaard and
Nietzsche to be the greatest ``worldview-psychologists,'' because they in deep and peculiar experience
have lived through the problems of existence, and because they in the first line emphasize the
genuineness of personal life, in whatever direction it then otherwise moves. And the two named thinkers
do indeed move in very opposite directions. --- One occupies oneself in Germany generally a great deal
with Kierkegaard, now that there has come a complete German translation of his works.\footnote{The
Spaniard Unamuno, in his remarkable book \textit{The Tragic Sense of Life} (English translation 1921), is
also influenced by Kierkegaard, whom he calls his ``spiritual brother.'' He reads him in Danish.} And it
is more his assertion that subjectivity is truth that interests, than his polemic against official
Christianity.

Jaspers's juxtaposition of Kierkegaard and Nietzsche gives me occasion to mention a lecture I gave in
Helsingfors in 1911, and which was later printed in the third volume of my \textit{Mindre Arbejder}. It
gives a comparison between Kierkegaard and Nietzsche. I first draw attention to the great opposition
between the two thinkers, but bring out that rivers flowing in different directions can well have their
source on the same mountain, and that philosophy is, after all, more interested in the sources than in
the course, more in the original life-experiences than in their interpretation. Common to Kierkegaard and
Nietzsche is that they both see far beyond life's ordinary, delimited, rational paths --- that they both
maintain the innermost and original in man as something that cannot be made rational, while it yet can
possess the highest power and value. They are symptoms of a need that stirs in the time, and they give it
a volcanic expression. I then seek to demonstrate a whole series of analogies between the two so different
life-thinkers, but cannot here go more closely into them. The treatise ends with the words I can still
hold fast: As long as life advances, its whole content will not be able to be wholly and entirely
arranged by thought, but it will constantly be a matter of keeping the mind open for the new that life and
experience can bring. --- Later I have made a similar investigation of the relation between Pascal and
Kierkegaard (in \textit{Tilskueren}, 1923). Here the analogies were closer, but the differences, due
partly to the opposition between Catholicism and Protestantism, partly to the different conditions of the
times at the appearance of the two great spirits, are therefore no less characteristic.

However much interest philosophy can and must have for individual attempts to take a position on the
problems of life, it yet belongs precisely to its tasks to guard the boundary of rational thinking, so
that Dogmatism shall not make its entry under the name of view of life or philosophy of life. But first
and foremost it must seek, with the means at its disposal, to become clear about the peculiar psychic
whole-formations that can step forth in the form of views of life.

\hnum{17} There are two different causes that attention has more and more directed itself toward
individual views of life.

The one cause one was already clear about in antiquity. It lies in all science's general character, which
Aristotle so strongly brought out. Science seeks understanding through the general laws of existence. It
seeks starting-points that make it possible to draw conclusions that hold for everything and everyone.
But life, and the whole of existence, reveals itself to us from scattered, mutually very different
starting-points, each conditioned by its peculiar relations. The single things or beings each appear
with their individual constitution, and general propositions can find application to them only
approximately and can in any case not exhaust their nature, not let all sides, all possibilities in them
come to their right. But what we call a man's view of life is now precisely, when it is worth anything,
something individual and personal, built up on its own experience and its own interpretation, acquired in
the struggle of life, through the work to reach the goal one has set oneself. Through experiences, now of
life's greatness, beauty, and happiness, now of its narrowness, ugliness, and suffering, there will
involuntarily form a picture whose form and whose color are determined by what one has seen, thought,
felt, and striven for. The scientific cognition that may be present can play a role too, but only as an
element woven into the picture, not as all-determining. In no man, not even in the greatest man of
science, is the view of life determined exclusively by science. In no one does humanity go up entirely
into scientific character.

There was a time when one had, as it were, views of life in stock. Their names always ended in -ism. One
had optimism and pessimism, idealism and materialism, monism and pluralism, individualism and socialism,
and so forth. And one supposed as a rule that one did not need many indications in order to know which of
the many drawers one should pull out to place a view at hand in the right spot. One assumed too that one
knew all possible directions for the human way of taking life. And yet it might be that life was richer
than that, and that the kinds of view of life found hitherto did not suffice.

There was a young man in Syracuse who gave himself the appearance of knowing Plato's conception in its
innermost essence. On this occasion Plato --- this great master, unsurpassed in the literary presentation
of philosophical thoughts --- stated that his innermost meaning could not be expressed in writing. Only
through long occupation with the thoughts, preferably in personal discussion together with their author,
could full understanding be reached, and a light may then suddenly emerge that can afterward burn by its
own power. --- More or less it goes, as Plato here describes, with all cognition, but it holds quite
especially of a man's innermost thoughts about life, as it has revealed itself to him.

\hnum{18} Besides the opposition into which science, by its general character, comes to stand to life,
another relation of opposition also makes itself more and more felt. The conduct of life demands
concentration on certain definite thoughts, gathering in one definite direction. But within science as
within spiritual life on the whole, the division of labor prevails more and more. There is in modern
culture no intellectual unifying power that can step up in the face of life. There will then be a need to
gather one's different experiences into an individual whole, even if this whole is not built up according
to the direction of science or of authority, but by what Kant has called ``the hidden art'' that works in
us without our needing to notice it. It is a building erected in the way that follows from the single
individual's nature and circumstances, and which is first known when the involuntary work of building has
advanced to a certain point. It is psychological and historical relations, not pure logic, that condition
this building's character. But experience shows that there is a certain kinship between consistency and
character, which makes the holding fast of a screaming contradiction between standpoints that seem to
proceed from experiences made, for most people, a spiritual suffering that can have lasting value only
from an ascetic position. A natural need will stir to get beyond the contradictions, so that what remains
of disharmony stands only as a witness to life's fullness and manifoldness. For many natures there
remains, indeed, a sting, a certain unrest that, after all, belongs to the essence of life. Giordano
Bruno brought out, in his work ``on the heroic affects,'' that there is an inner fire that is never
extinguished as long as one strives after a high goal. Heroic affect arises precisely when one does not
let oneself be deterred from striving toward a high goal because there is pain and danger bound up with
this striving.\footnote{See further on the remarkable writing of Bruno mentioned: \textit{Den nyere
Filosofis Historie}, 3rd ed., I, p.\ 135.} What the Renaissance's most typical thinker thus stated will
make itself felt in every view of life that is not built on too narrow a ground and does not fear coming
into the heights.

\hnum{19} Besides the general relations mentioned in the foregoing, which lead to individual view of life
coming to play a role, there lie in every personality grounds that bring it about that it comes to arrange
things for itself, even if science or religion gives no direction to it. They are grounds that also work in
the coming-to-be and development of science and religion, but which lie so deep in human nature that they
can be traced even where there is talk neither of science nor of religion. Some of these shall here be
drawn forth.

In the first rank must here be named the urge of expansion --- that is, the involuntary inclination to let
a striving, a mood, or a train of thought that one has once entered upon continue and unfold itself
further. The surplus of energy one disposes of is set upon what has thus caught the mind, and later
experiences are seen and interpreted in the light of them. It is the ``will's first thought'' of which
Dante speaks, and which can be so difficult to explain, because its origin is in the dim. Instead of
seeking and laying out wholly new standpoints, the way one has once struck out upon is followed. In science
too one provisionally sets out from the assumption that what has shown itself to hold in a series of cases
holds for other similar cases. It leads to the setting-up of a hypothesis, which is then tested more
closely by seeking out new similar cases or by drawing conclusions from the hypothesis and seeing whether
they agree with old and new experiences. It was an expansion that (whether the story of the apple's fall is
true or not) led Newton to attempt the application of the law of falling bodies to the moon's and then to
other heavenly bodies' motion.\footnote{See on this \textit{Den nyere Filosofis Historie}, 3rd ed., II,
pp.\ 203--205.} Science here asks ``Why not?'' like the child and the savage. But what sets science in
opposition to the involuntary unfolding of the life of thought in all of us is that it formulates the result
of its expansions in such a way that it can be the object of exact testing. Expansion can also make itself
felt in the domain of feeling-life, as when a deeply penetrating joy or sorrow comes to stand as a witness
of what one can at all expect from life, so that thereafter either optimism or pessimism becomes ruling.
According to the legend\footnote{Translated in Warren, \textit{Buddhism in Translations}, p.\ 56\,ff.} it
was the sight of the infirmity of old age, of illness and of death, that determined Buddha to withdraw from
the world. These sights were life-experience enough for him. For Pascal and Kierkegaard it was the story of
Jesus in Gethsemane that came to stand as the expression of what life could bring when one set oneself a
sufficiently high ideal. The Gethsemane mood became for them life's proper, true ground-mood, which had to
rule in all who would call themselves after Jesus.\footnote{Cf.\ my treatise \textit{Pascal og Kierkegaard}
(Tilskueren, 1923).} For other natures, like Spinoza and Leibniz, the fact that joy of cognition, joy in
intellectual work, can be attained may be enough to determine the whole mood toward life. In every serious
view of life there will be experiences that become decisive for the whole mood toward life, and if one
will understand a view of life, one must see to find what has thus become an ``ur-phenomenon'' for a man.

When Homer calls man a being that looks forward and backward, the relation is yet that we look forward
before we look backward. Life is lived in expectation before it is lived in remembrance. The surplus of
power with which life begins attributes from the first to representations that are called forth a practical
significance as announcements, and makes them into motives for action. It is first unfortunate, or at any
rate unexpected, effects of the actions that refer certain representations once and for all to the past, so
that no expectations are attached to them.

Expansion leads naturally to expression. The mood aroused demands expression. Georg Simmel, who in his
above-cited treatise explains the unrest of the time by the fact that the old forms are no longer a home
(\textit{Gehäuse}) for the urge that stirs, uses especially a certain direction in our time's art as an
example. According to his conception, so-called Futurism finds its psychological explanation in the fact
that ``the artist's inner movements are continued, just as they were experienced, in the work, or rather as
the work itself.'' It is not to be cast into any form, and no imitation is to work along. (\textit{Der
Konflikt der modernen Kultur}, p.\ 111.) This consideration can be extended. Where the unfolding of life
takes place involuntarily, it yearns after expression before the experiences are discussed through
reflection. The simplest judgments are exclamations,\footnote{\textit{Den menneskelige Tanke}, p.\ 84.} and
there are linguists who espouse (see the last chapter in Otto Jespersen's \textit{Language}) the hypothesis
that the first stage of linguistic development is poetry, in that it is predominantly emotionally
conditioned. Just as now the closer formation of the judgment takes place through the necessary
specialization and the drawing-forth of conditions, so it depends on whether the view determined by
expansion from a significant experience can keep itself up in the course of life. There is here a point that
forms an analogy to the confirmation Newton's expansion received from the fact that continued experience
agreed with what deduction from his hypothesis demanded. It is naturally a distant analogy, and precisely
this distance from the exact is characteristic of the conditions of views of life. Left to itself --- that
is, without new expansions from experiences similar to the founding one --- the movement aroused in the
mind will be able to last until the powers are exhausted. But there can also come new expansions from
experiences of a wholly different kind than the earlier, and there will then be able to arise dividedness,
doubt, and direct strife in the mind. It is the time of problems. It is here as everywhere discontinuity
that occasions problems, while continuity is the condition for rest and security in the once-trodden ways.
Henri Poincaré has stated that in scientific discussion the struggle between discontinuity and continuity
will last as long as man will think.\footnote{\textit{Les conceptions nouvelles de la matière}, p.\ 67 (in
the collective work \textit{Le matérialisme contemporain}).}

The same holds in the case of views of life. One should be ``true to oneself and to one's best striving,''
says a poet. But when is our self finished and mature, so that it is merely a matter of fidelity to it as it
has become? And which striving is the best? Life can constantly bring us into new situations in which our
personality is to stand a new test. As Jaspers\footnote{\textit{Psychologie der Weltanschauungen}, 2nd ed.,
p.\ 230.} has beautifully shown, it is especially in ``boundary situations,'' where the hitherto experienced
and fixed values are threatened with destruction, that the personality is to stand its test. Such ``boundary
situations'' are struggle, death, chaos, guilt. To these situations suffering is attached, and it then
depends on whether there is energy enough to stand firm despite it. The detailed depiction Jaspers gives of
this whole question is plainly due to the study of figures such as Kierkegaard and Nietzsche. (Cf.\ above,
\S\,16.)

Kierkegaard has set up a measure for the height of spiritual life, which consists in the energy with which
the oppositions life can present are held fast, even if they should be so sharp that life becomes a
suffering. At the foundation here, in Kierkegaard, lies the concept of synthesis stemming from Kant, and
there is then set up a graded series of synthetic energies, right up to the energy that is able to conquer
even if life's most violent oppositions collide with one another.\footnote{See \textit{Søren Kierkegaard
som Filosof}, 2nd ed., pp.\ 119--126.} By the setting-up of this measure Kierkegaard has given an important
contribution to the psychology of views of life, even if his application of the measure is stamped by the
Gethsemane mood which, for him as for Pascal, designated the high point of spiritual life, but which
presupposes the rejection of all tasks that otherwise can set the mind in motion in human life. If it is so
that the highest entered the world at one definite point of time and at one definite place, and that there,
in the decisive moment, it was denied, even by those nearest, then the Gethsemane mood is the only worthy
mood toward life. Only one must then not content oneself with seeking to call it forth once a year, but it
must be ruling constantly and determine the relation to everything else that could give human life value.

Taken in itself, the Kierkegaardian measure retains its great significance, even if it does not beforehand
exclude strife about what stands highest on the scale of value.\footnote{Such a strife was once carried on
by Baggesen and Grundtvig. They disagreed about what was the highest, especially about the place of
cognition in spiritual life, but they lacked a measure. Cf.\ \textit{Danske Filosofer}, pp.\ 66--72.} It is,
in spiritual strife, in the end the only standpoint one has to hold to. The energy with which life is
conducted --- the fullness of the content that can be worked together, and thereby of the oppositions that
can be lived through without life being burst apart --- it is in the end what it depends on. But it will
always be a defect, and in the end a misfortune, when an essential part of life's content must be denied in
order for one to secure what one calls the one thing needful. From a human standpoint the one thing needful
can precisely only be the energy with which one takes life, in that one, according to one's capacity, seeks
a form in which its content can come to its right in the single personality.

There is one thing the indicated measure necessarily demands; it is what the movement in the direction of
the principle of personality, introduced by Carlyle, requires: honesty, personal truth. The point is not to
mistake one's place on life's scale, to attribute to oneself a relation to ideals one does not in spirit and
truth espouse. It is naturally a difficult task that is hereby set, but it must be taken up. By a single
example it can be illuminated, when we think of the opposition between tragedy and humor. The Gethsemane mood
is tragic. He who at the very first hand lived through it stood firm and walked the way of suffering to the
end. But it stands with this mood as with all tragedy. In an earlier work I have, as opposed to it, described
a way of taking life that gives the great and valuable the highest position without overlooking the
sufferings and imperfections, and which in its inward absorption in the great can even have a smile at its
disposal in the face of the small. The ground-mood I have thus sought to characterize I by no means regard
without further ado as the ``right'' or as the ``best,'' let alone as the only possible, or as the one in
which all life's elements have come to their right. I maintain, and that not merely as a concession, that a
tragic conception of life is constantly possible and can be unavoidable. It is, after all, possible that a
life cannot be harmonized without danger of narrowing. The humorist yields in the end to the tragic hero who
does not fail his task, even if all harmonies are burst apart. There is nothing to prevent the humorist
himself from becoming a tragic hero and standing the highest test to which men can be put. But it is of
decisive significance to hold fast\footnote{Cf.\ \textit{Den store Humor}, pp.\ 104--108.} that the tragic
shall come ``of itself,'' by the course of things --- the course of things in which man, in order to be able
to solve his own task, has had to intervene. It is a fate that comes, but it is not something that it should
be everyone's duty to experience.

In the interesting characterization Gorky has given of Leo Tolstoy, he says that what especially repelled
him in the great writer was his desire to come to suffer for his ideas, in order that these might thereby
exercise greater influence.\footnote{\textit{Reminiscences of Leo Nicolayewitch Tolstoi. By Maxim Gorky}
(London, 1920), p.\ 40.} For it is one thing that a tragic fate can be foreseen, another to wish for it when
the necessary course of things itself does not bring it. It comes like an earthquake that cannot be called
forth. A tragic type can therefore, just as little as a humoristic type, without further ado be set up as
the high point on the scale. Neither tragedy nor humor are mere consequences of psychic energy displayed,
but their possibility rests for a great part on relations that do not stand in our power.

\hnum{20} Both in the face of the view of life that sets tragedy, and that which sets humor, as the highest
on the scale, life can step up in its inexhaustible manifoldness. Is it possible to comprehend this richness
in one single yearning, or in any harmonious ordering whatever? There can constantly emerge new sides to
life, and even the most contemplative consideration, the greatest life-experience, cannot get everything
exhausted. There will constantly subsist a conflict between the urge to gathering and the urge to widen the
horizon. The concept of the whole points everywhere, where it occurs, beyond itself. It is a boundary
concept, as I have sought to show in the treatise on Totality as a Category (Ch.\ 6). This holds also of the
whole that personal life strives to become. And yet it is with right that the measure laid on the different
views of life sets the fullness of the content as an essential determination. It is easy enough to
comprehend a meager content.

While expansion essentially leads to running the line out in a definite direction, the other side of the
measure demands that account be taken of the many different lines of development that can stir in the mind as
more or less restless tendencies. But the work of life is now, like all work, subject to the law of division.
Intellectual development sets its demands, as does the aesthetic and the religious, further the economic
securing of life and participation in civil life. Each of these interests can demand a whole man, but the
single ones should well be able to be whole men without denying any of these interests, which are not always
easy to bring into mutual coherence. It is a particularly happy case where the task directly corresponds to
the energy at one's disposal; but it then depends on whether the task is also set sufficiently high. A happy
case it can also be where there is far more energy at one's disposal than the task demands; here too it then
depends on whether one has taken all considerations into account when the task was set. A tragic case it
becomes, on the other hand, when the energy falls far short of what the task demands; it can presage
spiritual or bodily destruction. Thus the urge to set oneself comprehensive ends (one might call it the urge
of complexity) steps into a certain opposition to the urge of expansion.

Also from this side I think Simmel was right when he found, as the essential ground of the time's spiritual
unrest, that the old frames, the old spiritual homes (\textit{Gehäuse}), no longer suffice. Here Simmel
distinguishes himself from the pessimists who deny our time all culture or trace dissolution everywhere. He
sees no ground why power, striving, and interest should not be present --- the same powers that produced the
old culture in its different forms. To this it is, I think, to be added that there is indeed much outdated in
the old culture, but there are besides many elements in it that have not been taken up and experienced in an
independent way. How comparatively slight is the role the world's unsurveyability has played for views of
life! One scoffs at antiquity and the Middle Ages for their limited horizon, and yet most people hold to the
old mythological representations of heaven and earth. And in vain have great spirits proclaimed the gospel of
eternal striving; one relies in the end on de-souled dogmatic representations. Furthermore, absorption in the
works of the great artists and thinkers will constantly be able to breed new motives, even if the tasks the
ever-advancing life sets must stand in the first line. Here then step up tasks that are not in themselves in
conflict with each other, but which in life often step up in hostility toward each other. It is these to
which Goethe has given expression in the two lines:

\begin{quote}
\textit{Aeltestes, bewahrt mit Treue, / freundlich aufgefasstes Neue!}
\end{quote}

[Oldest, preserved with fidelity; new, kindly received!] Between fidelity and openness there is, after all, a
great struggle in the world. While the urge of expansion leads most nearly to the continuation of the once-given
direction, the urge to fullness of content leads to openness in the face of new directions.

\hnum{21} In a certain opposition both to expansion and to complexity stands the repetition that life
unavoidably brings. The urge of expansion can sometimes have enough in one mighty experience, in whose light
everything later is seen. But if the energy is not great, repetition of expansion's expression can become
weaker. On the other hand (as already mentioned), repetition of the original experience can be beneficial,
because very disparate experiences will weaken the expansion. For the urge of complexity, which sets out to
widen and supplement the content won hitherto, repetition of the same movement will be able to be unfortunate.

For every view of life a constant direction of the life of thought, feeling, and drive is a characteristic
quality. The direction must be identical, although expansion runs through different stages, and although the
urge of complexity constantly seeks new content. Against this it does not conflict that there are views of life
that make precisely their principle that repetition is to be avoided at any price and every moment to have a new
experience. For this cannot be reached without a virtuosity that demands practice and so far repetition too.
This is seen both from what is reported about Aristippus and from Kierkegaard's depiction of ``the aesthetic
stage.'' Art and practice are required in order that remembrance and expectation shall not hinder full
absorption in what the moment brings.\footnote{A comparison between Aristippus's teaching and Kierkegaard's
``aesthetic stage'' I have made in \textit{Søren Kierkegaard som Filosof}, 2nd ed., p.\ 84.} It is only the
young mind that can say, as Goethe in the Marienbad Elegy lets the young woman say:

\begin{quote}
\textit{Stund' um Stunde / wird uns das Leben freundlich dargeboten. / Das gestrige liess uns geringe Kunde, /
das morgende, zu wissen ist's verboten.}
\end{quote}

[Hour by hour life is kindly offered us. Yesterday left us scant knowledge; the morrow, to know it is forbidden.]
Later in life both the past and the future announce themselves, and energy is required to keep them at bay in
order to meet the moment's demand. ``\textit{Du hast gut reden}'' [It is easy for you to talk], the old poet
therefore answers to the words he has put in the young woman's mouth.

It is characteristic that both Kierkegaard and Nietzsche saw, in the repetition that on account of life's
conditions cannot be avoided, a great danger for their views of life. Kierkegaard, to be sure, praises it when
it is met by a will to take up life anew under the same conditions, or when one hopes with religious trust; but
there is then required a commitment which, according to him, cannot be explained by psychological principles.
Nietzsche saw in repetition the great danger for the affirmation of life. He feared leveling, boredom, and
dullness as its effects, and his Zarathustra dies first reassured when men have promised to will life despite
its constant cycle. In opposition to these modern modes of view, the philosophical emperor Marcus Aurelius found
precisely a reassurance in the fact that everything repeats itself. Partly one can then live without fear of
meeting wholly unforeseeable hindrances on one's way; partly one can feel oneself one with the power that works
in nature's rhythms, and there can develop a disposition in which one says to this power: everything your seasons
bring is fruit for me.\footnote{Cf.\ my treatise on Marcus Aurelius in the third volume of \textit{Mindre
Arbejder}. In the same volume is found the comparison between Kierkegaard and Nietzsche.}

Every view of life must have an intellectual background, this being more or less formed and conscious. There
must have formed a whole-picture of existence, as this is, in greater or smaller compass, accessible to the
individual. It follows already from the fact that everyone must be able to orient himself in the world he lives
in, and must be able to distinguish between dream and reality. This is possible only when the same thing occurs
under the same relations, and it is experienced only when the same relations recur. From the intellectual side
every view of life is therefore conditioned by repetition.\footnote{Cf.\ the section on intellectual culture in
my \textit{Etik} [Ethics].}

For feeling-life the matter presents itself somewhat differently. The vehemence of feeling can be weakened by
repetition of the occasions to which it is due. On the other hand, inwardness can grow through repetition, when
the object of feeling only has a rich content that can unfold itself successively. When life is lived in work for
great tasks, interest can rise at the repetition of the same posing of the problem. Repeated reading of a
significant work, or repeated viewing of beautiful nature or art, can increase the joy. There is won, to use
Goethe's expression, in inner growth what is lost in voluptuousness. But here can lie the germ of a different
view of life, according as the intellectual or the emotional element has the preponderance, and according as,
within feeling-life, vehemence or inwardness has the greatest value.\footnote{This consideration I have already
set forth in my \textit{Psykologi} (abridged ed., pp.\ 188--191).}

As a condition for recognition, repetition easily leads to the new that is experienced being interpreted from the
old that has been experienced. But for an advancing view of life it can be of great significance that, conversely,
the old is interpreted from the new. Here too lies the possibility of an opposition between different views of
life.

\hnum{22} In every view of life it will show itself that the three moments drawn forth above --- expansion,
complexity, and repetition --- work along. But each of these moments can, by being predominant, determine a view
of life of its own kind. In Plato's depiction of Eros (in the Symposium and in the Phaedrus) expansion, which is
driven forth to ever higher stages, is thus ruling, whereas in his books on ``the State'' he sets up a type of
complexity that leads to harmony between different spiritual tendencies. In more recent time Schiller represents
the type of expansion, Goethe the type of complexity or harmony.\footnote{As Chr.\ Sarauw shows in his beautiful
and substantial book \textit{Goethes Augen} (Vidensk.\ Selsk.\ Historisk-filologiske Meddelelser, 1919), pp.\
134--139, Goethe occupied himself his whole life with the opposition between complexity and self-limitation. In
his youth the urge to embrace a world, later the urge to self-limitation, had the preponderance. --- What is here
called self-limitation belongs most nearly under the meaning attached above to repetition in relation to the view
of life.} When repetition plays the greatest role, and the whole mind reassures itself by it, there develops
either a Stoicism such as Marcus Aurelius's or a philistinism such as can be met frequently enough.

But common to all views of life is that innermost there rules an interest in one or more sides of life. Thereby is
determined the direction in which the unfolding goes, the supplementing is sought, or repetitions are desired. It
need not be the individual's interest as isolated that works; it can be an interest aroused through living oneself
into the tribe's or the estate's tradition, that determines the direction, as Durkheim has shown in the case of
the totem-worshippers.

The subjective character of every view of life thus rests not merely on the fact that the starting-point for
expansion, supplementing, or repetition can be different in different individuals, but also on the fact that it is
a value, an interest, that determines the very starting-point. Each individual and each circle of individuals is a
little whole whose conditions of life condition particular values. In the practical synthesis that every view of
life is, many elements thus work together. And the analysis of what has thus, by life's involuntary ways, joined
itself together can often be very difficult to carry through.

Between the different views of life stands the great struggle in the world. The one man here often does not
understand the other at all. Into a foreign view of life one must first and foremost \textit{live} oneself.
Greatest is the difficulty in the relation between different religions, which each bear the stamp of a peculiar
culture and build on particular traditions. But within the same religion too difficulties can arise on account of
the confessors' individual experiences and characters. The living sense for personal peculiarity in its connection
with conditions of life and the course of life is the noblest flower of spiritual culture. Under the presupposition
that such a living sense is present in the psychological researcher, he will then go over to analyzing the whole,
the peculiar crystallization, that he has before him.

The first standpoint in the strife between different views of life will be the one that concerns genuineness, the
personal truth in each view of life for itself. The point is whether the spiritual building has come to be in such
a way that the single elements have joined together in a natural manner, so that self-deception, on account of a
desire to make a figure in one's own or others' eyes, has not worked along. This is naturally a very difficult
matter to judge. There is here presupposed a sense for the interplay between the conscious and the unconscious and
between the arbitrarily assumed and the involuntarily arisen. And the difficulty is here so much the greater, since
there must, within every view of life, be a distance between the goal one has set oneself and what has hitherto
been attained. Likewise there will be able to be new thoughts and modes of view that have not yet grown completely
together with the real self, because it requires time before they can be wholly appropriated. The point will be to
become clear about the whole direction in which the personal striving goes. The influence that Søren Kierkegaard,
and later Ibsen's \textit{Brand},\footnote{See the interesting demonstration of this work's influence on Strindberg
in Martin Lamm, \textit{Strindbergs Dramer}, I, 44--49.} had on the generations that stood nearest to them, had
provisionally to remain indefinite, merely disquieting and inciting. Only gradually was it possible that the new
thoughts in single individuals could be worked in as links in their striving, often perhaps in wholly other forms
than the great thinker and the great poet had meant. And there were no doubt many who consoled themselves that the
time had not yet come to apply these thoughts. In others perhaps the unconditioned demand of genuineness led to
nihilism, because self-searching could not be brought to a conclusion.\footnote{Cf.\ on this Jaspers,
\textit{Psychologie der Weltanschauungen}, p.\ 288\,f.} The matter is that it holds for genuineness as for every
other ideal, that the beginning of applying it must be made in a single domain, preferably the nearest, and that
other domains therefore do not provisionally get its stamp.

Another standpoint that must here be taken is the necessity of supplementing. The more individual and genuine a
view of life is, the sooner must the insight arise that not all can apprehend life in the same way. As no work of
art teaches us to know nature, especially human nature, completely, so it stands too with the work of art that
every view of life is to be. Only the whole race's life and history can give us an all-sided acquaintance with the
human. The single individual will be able to get valuable motives for self-development by gaining insight into how
others, from the conditions given to them, have developed. Such insight will now warn, now tempt, now lead to new
possibilities in one's own self being discovered. But the single individual's personal view of life will then,
when it has developed in a natural and honest manner, be able to stand as a link in the race's unconscious and
conscious work to develop a life of personality in which all the strings of life's experiences are
touched.\footnote{Cf.\ my \textit{Religionsfilosofi}, 3rd ed., p.\ 192.} As the single individual, despite his own
idiosyncrasy, can learn from others' idiosyncrasies, so perhaps his idiosyncrasy can directly or indirectly come to
the good of others. He can, as it is said in Longfellow's ``Psalm of Life,'' leave behind him ``footprints on the
sands of time'' which perhaps later another, who also wanders about with his idiosyncrasy, can discover and gain
courage from.

The modern theory of relativity, which popularly is often misunderstood as skepticism, leads precisely to seeing
the necessity that single individuals, under certain relations, cannot have the same conception.\footnote{See
\textit{Den menneskelige Tanke}, p.\ 304. --- \textit{Relation som Kategori}, p.\ 75\,f.} The principle of
personality leads in its domain to an analogous result: men, after all, do not all have the same lesson to learn.
Each single one has his own. The principle of personality contains precisely the truth that everyone must approach
truth by his own way; but it leads at the same time to an ever better understanding of these different ways. Each
single one is led to regard himself from a universal standpoint, as one of the many ``monads'' of which humanity, after
all, consists.

\hnum{23} The more genuine, the more involuntarily built a view of life is, the sharper can the collision become
when such a life-synthesis meets science's syntheses built on logical analysis. But such a meeting can be
unavoidable, and it is no good sign of the time, but a lack of intellectual honesty, when one acts as though there
were here no problem at all. The immediate security with which life begins, and which is easily preserved as long as
the horizon is narrow, cannot always subsist when it is set in the face of the great intellectual development that
is, after all, a necessary element in human culture. In the popular discussions of this subject, he who in stillness
follows them must often be frightened at the lack of understanding of what the problem is about, that can appear. One
asserts from the one side that science has once and for all settled the problems belonging here, and from the other
side that science means nothing at all in the struggle of views of life.

Against this last assertion it must first and foremost be remarked that the very endeavor to win scientific insight
can stand as that which gives life worth, in that it ennobles the drives and directs the mind's yearning toward one
goal. It is in any case a type of view of life beside other types. Expansion here takes its starting-point from the
experience of the joy of cognition, and through the advancing work that cognition demands, life becomes constantly
new. Despite all disappointments and all inner and outer resistance, even the most modest participation in, or
appropriation of, men's great intellectual common work can breed trust and hope. Now there are indeed only few whose
view of life can predominantly be built on this foundation. But Carlyle has rightly said that it is not necessary
that a man should himself have discovered the truth toward which his faith is directed; it is not novelty, but
sincerity, in which true originality consists. (``The merit of originality is not novelty, it is sincerity.'')
Intellectual honesty --- that is, penetrating reflection on the significance of the insight the single individual
can procure under his circumstances --- must yet be an element in every serious view of life, and it cannot be
decided beforehand how great an influence such reflection can get on the view of life. One cannot, like Kierkegaard,
distinguish once and for all between ``essential'' cognition --- that is, the cognition that is of significance for
personal view of life --- and ``unessential'' cognition (that is, the cognition that has a purely objective
character). That, for example, the great widening the world-picture has received through the science of more recent
time, from Copernicanism to the theory of relativity and the doctrine of electrons, should have no significance for
men's way of taking life, is incredible. Pascal, who himself was, moreover, taken with the widening the world-picture
had received (although he was not a Copernican), advised against absorbing oneself too much in the sciences; he
thought of writing against those ``\textit{qui approfondissent les sciences}'' [who go deeply into the sciences]
(fr.\ 76, éd.\ Brunschvicg). Grundtvig advised against occupying oneself with Copernican thoughts. On the other
hand, the widened world-picture in the Renaissance's thinkers developed a feeling of the exalted, not merely on
account of the fact that the world one lived in was so great, but also on account of the fact that human thought had
discovered this greatness. High-mindedness became a disposition that for the thinkers of that day was called forth by
scientific work. (Telesio: \textit{sublimitas}; Descartes: \textit{generositas}.)

One might ask whether the time must not come when there can no longer be talk of views of life, because science can
settle all questions, even the most personal we can raise. According to what has been developed in an earlier section
of this treatise, this question must be answered with no. Science, both the natural-scientific and the humanistic,
will indeed at all times come to stand in the face of individually formed views of life, peculiar psychic wholes
that have arisen predominantly involuntarily, through a spiritual natural process. Their existence as such science
must provisionally acknowledge. How much sympathy the single researcher will feel in the face of certain definite
wholes of this kind will rest on causes that need have nothing to do with science. But it becomes a task for science
here --- as in the face of the organism and the personality --- to seek to explain to itself how such a whole has
come to be. It is a psychological-historical task that the more recent research in literature and religion has
acknowledged, and which demands the cooperation of many qualities. From the side of life there will often (and that
not merely where it is a matter of religious views of life) be looked with hostile eyes on such
psychological-historical investigations. But there lies here an irrefutable task. There appears sometimes a so-called
``philosophy of life'' that can be just as dogmatic as the old orthodoxy and far more dogmatic than the science it
would combat. And its tone is often of the kind that got Kant in his old days to write a treatise ``\textit{Von einem
neuerdings erhobenen vornehmen Ton in der Philosophie}'' (1796) [On a recently raised superior tone in philosophy].
Kant is malicious enough to recall that one often considers it ``superior'' not to work, and he supposes that it fits
well those who strike the superior tone in philosophy. He speaks especially of those who would philosophize from their
feeling and look down on thorough and objective thinking. But other ways too can be gone. One tries to deny the
existence of decisive problems and attempts to go back to the problem-free state that is so natural for a childlike
consciousness. Or else one asserts that these problems were solved long ago. Or one constructs in a trice a whole
little metaphysics that is in all comfort to give the solution, but which often is only a web of more or less poetic
analogies and false conclusions.

\hnum{24} The individual character of the view of life by no means sets any absolute boundary for science. Even if
science cannot exhaust personal life, it can yet find facts and standpoints that illuminate it. It can perhaps find
the leading motives and the means of thought that were at one's disposal at the given time. And it can investigate
the relation between the individual experiences and interpretations on the one side, currents in time and society on
the other. This kind of research is still in its infancy, partly because it often lacks sources, partly because
dogmatic prejudices veil the gaze, partly finally because artistic capacity is needed to give a vivid picture of a
striving personality in its relation to the tasks and riddles of life.

The question has been raised whether epistemology can take a position in the face of the ``philosophy of life'' that
especially in Germany makes itself heard, and which indeed often has a dilettantish character and rests on a fashion,
but which is yet ``a sign of a deep longing in our time,'' an attempt at ``new voyage over the infinite sea whose
waves break against the little island where secure, universally valid, logically formulated cognition dwells and
builds.''\footnote{Frischeisen-Köhler in \textit{Kant-Studien}, 1921, p.\ 135\,f.} --- If I am to answer this
question, I will state that philosophy (especially epistemology) stands in the face of the involuntarily built views
of life as it stands in the face of science. In both places it asks after spiritual powers that are at work, and after
the relations under which the work is done, rather than after the results that have been reached, and after the
definite forms in which these results are at a given time set forth. In the face of science it asks about the
presuppositions that lie at the foundation and about the ways along which it seeks to reach its goal. In the face of
views of life it interests itself more in the need that gets expression through them than in the clothing in which they
appear. In both respects the inner powers will keep themselves longer through the ages than the forms in which they are
at a given time encapsulated, and which are often more solid than they need be. Also in outer nature it is not the
forms, but the energy, that subsists. It is in this spirit that philosophy meets a personal view of life, regarding it
as a whole proceeded from life, whose bearing powers it is a matter of finding.

This is not to be understood as though philosophy stood raised above what leads men to form views of life. Philosophy
is --- in the philosophers themselves as in others who come into contact with it --- only an element in the single
individual's spiritual housekeeping. As in every researcher, there will also in the philosopher stir other spiritual
needs than the purely intellectual, however great a preponderance these last may have in the mind. A researcher's
personal life-experiences can lead him in a different direction than his work of thought would in itself lead. Only a
few of the figures the history of thinking presents to us yield a picture of full harmony between the life of thought
and the other sides of spiritual life.\footnote{Cf.\ for a single point \textit{Den store Humor}, Ch.\ 9.} Of a
complete fusion there can here scarcely be talk in any researcher, but there can be such an organization of powers and
interests that the personality's life can preserve its whole-stamp, even if there is one interest that has the
preponderance and directly or indirectly determines the place other interests get within the whole personal
existence.\footnote{On the difference between fusion and organization see \textit{Den store Humor}, pp.\ 26--34.}

\hnum{25} It is, as earlier remarked, of great significance for the science of spirit that personal life unfolds
itself powerfully and freely. Thereby, and only thereby, does psychology get a rich material. But how does it stand
with the science of spirit called ethics, which strives after a rational grounding of rules for the conduct of life and
for action? Will the difference of views of life not exclude that such an attempt can succeed, since there will be just
as many ethics as views of life?

The difficulties that here appear, and which perhaps will constantly hinder the carrying-through of a strictly rational
ethics, hang at their ground together with the fact that the concept of value, with which every ethics must operate,
points back to the concept of the whole, in that the ground of value in each single case lies in the relation between
a whole (an organism, a personality, a society) and the condition for its subsistence and development. (Cf.\ above,
\S\,14.) Value is acquired by everything that brings the elements of the whole into a relation to one another favorable
for its continued subsistence and development. But since wholes are of different nature and their subsistence and
development set different demands, the values too will be different. The struggle for life between organisms, between
personalities, and between societies presupposes oppositions in these wholes' conditions of existence and thereby in
the values they maintain. From each single whole's standpoint a consistently carried-through scale of value can be
formed, which does not necessarily agree with the scales of value that must be formed from other wholes, especially not
with the groundings that lead to these scales of value.

It is not merely on history's great stage that the struggle between values takes place. Within each single whole there
can again appear smaller wholes, hence new value-centers: organs within the organism, drives and endowments within the
personality, strata and estates within society. Here too value can stand against value. A single example we had above in
the relation between intellectual value and other individual values. Just as the division of labor entails
difficulties for social ethics, so the personality-life's different interests entail difficulties for individual ethics.
A peculiar difficulty arises when the ground-value --- that is, the value that within a whole conditions all other
values' place on the scale --- is about to vanish, and a new ground-value has not yet fixed itself sufficiently to take
its place as decisive for the rank-ordering of the values. The transition to a new ground-value's dominion often takes
place unconsciously, or at any rate involuntarily. But then there can, when reflection stirs, be doubt about how far the
development of the new ground-value has advanced, and it will in any case often be impossible to discover beforehand
whither the development leads. The larva that crawls in the dust cannot, as Goethe says, imagine what it is to be a
butterfly. The higher stage is, after all, from the first not in existence, and only in faith and intimation can we feel
ourselves assured that we are on the way to it. But how does it then go with the rational norm? Is it to be built on the
old ground-value, which is vanishing, or on the new, which does not yet exist? --- Still more sharply does the problem
appear when, for the individual's own consciousness, there is a hostile opposition between the old and the new
ground-value, and when the pain of leaving the former cannot provisionally be outweighed by the security and joy the new
can occasion. Then there rules a mood of homelessness, which can be fateful.

Such difficulties appear now not merely in attempts to set up a rational ethics, but can also come forth for every
serious view of life, and everyone must here struggle on his own account; only indirectly can help be had from without.
What philosophy can here render will first be to point to what the history of ethics sufficiently establishes,\footnote{Cf.\
\textit{Platon's Bøger om Staten}, p.\ 29.} that it is often only through opposition, struggle, and doubt that one can
reach the ground-values, and thereby the fundamental thoughts, that can be held fast as normative. Next, philosophy will
draw attention to the fact that there is a comparative ethics that not merely characterizes the different norms that have
in fact been set up, but seeks a measure for judging between them. The possibility of such a measure is already given by
the fact that no view of life and no ethics can avoid striving after consistent carrying-through. The relation between
ground-value and the single, more special values must satisfy the demand of logic as well as that of psychology. Directly
or indirectly the ground-value is to determine the single actions, if an inner inconsistency, which in the long run will
become unbearable, is not to be the consequence. Such an inconsistency Goethe describes when he says of his young friend
(Karl August): ``\textit{Noch ist, bei tiefer Neigung für das Wahre, ihm Irrtum eine Leidenschaft}'' [Still, with a deep
inclination for the true, error is to him a passion]. That means, after the manner of expression used in the foregoing,
that there is conflict between the ground-value and the valuations that are followed in certain special cases. It is
naturally not enough to point out the inconsistency, but this can breed a sting that can help to reach greater harmony.
--- Finally, philosophy will be able to refer to the fact that it is very well possible that the single moment's, the
single urge's satisfaction in the individual, or the single individual's welfare within society, can be fully attained,
since precisely thereby the individual or the social whole can get conditions for subsistence and development. Through
this possibility the individual's ethics can, despite its relative independence, be subordinated as a link in social
ethics. By this way one comes to the idea of a realm of humanity, a realm of free personalities whose essence is developed
through the work for the larger whole, while from the latter's side the work is arranged in such a way that a harmonious
development of the single personalities becomes possible. In the attempt to present an ethics on the foundation here
indicated, one must yet be conscious that it is only one of the in-themselves-possible standpoints that has been laid as a
foundation. But if the attempt is made with clarity, it will be able to make the conflict between the different standpoints
more fruitful.\footnote{See further the third chapter in my \textit{Etik}.}

\chapter{The View of the World}
\label{ch:5}

\begin{center}
\textit{(Metaphysics; Cosmology.)}
\end{center}

\hnum{26} However involuntarily and purely personally motivated a view of life has come to be,
apparently independent of systems and dogmas, and however much it can be combined with a consciousness
of its purely individual character, experience yet shows that it easily acquires the character of a view
of the world, a cosmology, or, to use the current word, a metaphysics.

It is now fairly generally recognized that metaphysics, when we take the word in its strictest sense, as
the doctrine of the innermost essence of existence, is not and cannot be science. Philosophers have seen
this before other cultivators of science, but philosophy has not always been able to hold fast this
insight. An involuntary tendency stirs to see the whole essence of existence, of reality, revealed in a
single part or side of it. Existence as a whole cannot be a subject for our cognition; there is no
standpoint from which it can be surveyed in its wholeness, and we do not even know whether all its parts
and sides are represented in our experience, nor whether it really is a whole. Our research constantly
works with different subjects that cannot be so reduced that a single subject could underlie all the
others. Through the lawfulness we find, existence does indeed appear as a unity, but the manifoldness of
the subjects cannot be derived from it. Besides this opposition between unity and manifoldness, there
appear in our experience the opposition between spirit and matter and the opposition between subsistence
and development as problem-setting. In my book \textit{Den menneskelige Tanke} I have discussed the
problems belonging here and on the whole given a fuller account of my conception of the metaphysical
problem than I can here. That there is and remains one or more problems here, I constantly hold fast. The
three relations of opposition named are constantly in conflict in our research in the special domains; but
metaphysics first arises when one draws the opposed terms out of the special connection and tries either
to derive the one term from the other or to find a third from which both terms could be derived.

In earlier times one was often inclined to call all philosophy metaphysics, and this usage is sometimes
still applied. But there has, in the philosophical domain as in other scientific domains, taken place a
division of labor that has led to the development of epistemology, psychology, and ethics as independent
disciplines, and the metaphysical problem has thereby been separated from other philosophical problems,
although it cannot therefore be said to have fallen away.

All our thinking sets out from a subject as given or presupposed and then seeks to find predicates that can
give a closer determination of it. But since now existence itself cannot lie before us as a subject, not
even as a presupposed subject, the whole weight comes to lie on the predicate that could be assumed to
express its essence. The judgments that can here be formed are pure predicate-judgments without a proper
subject. It is a part or a side of existence that for thinkers comes to take so prominent a place that it is
assumed to stand as the expression of the whole of existence. One here works with an analogy, in that one
supposes that a part or a side of existence stands related to existence as a whole as, within our
experience, a part or a quality stands related to a thing. But existence as a whole is, after all, no
``thing,'' when by a thing one understands a subject given or presupposed in our experience. A presupposed
subject existence as a whole cannot be, since we cannot form any rational concept of it, such as of the
``things'' for which we in science seek new predicates.\footnote{See further \textit{Bemærkninger om den
platoniske Dialog Parmenides} (Vid.\ Selsk.\ Filosofiske Meddelelser, 1920), p.\ 56\,f.}

When we nonetheless attribute to existence certain qualities, we treat it in analogy with subjects that
occur in our experience or for our research. Leibniz already maintained that all metaphysics builds on
analogy; only he did not see that he himself, by his treatment of the concept of substance (see above,
\S\,2), had eliminated the foundation on which all earlier metaphysics had built, in that it set out from
the assumption that there was something ``in-itself-existing,'' and that it was now merely a matter of
finding the predicate to it.

Analogy plays a great role within science. There is here to be distinguished between exact analogy or
proportionality, which belongs in arithmetic, and rational analogy, which finds application when we find in
experience that to every link in a series of subjects there corresponds a definite link in another series
of subjects, and that this relation does not change however far the two series are continued. In my treatise
on the concept of analogy I have adduced as examples: the relation between arithmetic and geometry, the
relation between formal and real science, and the relation between natural science and the science of
spirit. Such rational analogies are at once hypotheses and methods of work. But they all presuppose that we
can form two series, both of which can be continued, each according to its own law and yet in a
corresponding manner. In the metaphysical problem the difficulty is the great one that we lack the one
series. Here lies the boundary between science and metaphysics. The predicates one supposes oneself able to
attribute to existence as a whole really express only the significance the predicate-concept in question has
won for us on theoretical or practical grounds. If practical grounds are predominant, the metaphysical
problem passes over into the religious problem.

In my treatise \textit{Filosofiske Problemer} [Philosophical Problems] (University program,
1902)\footnote{More fully I have later characterized the four problems in \textit{Den menneskelige Tanke}
(1910).} I have nonetheless placed the metaphysical problem on a level with the epistemological, the
psychological, and the ethical problem, in that I supposed that all thinking, whatever problem it discusses,
strives to carry through a continuity, so that the urge to continuity grounds an analogy between the
problems with which philosophy has to do.\footnote{A related conception has been set forth by Dilthey,
\textit{Wesen der Philosophie} (Kultur der Gegenwart, VI, I, 1907), p.\ 34\,ff., where it is maintained
that within all principal problems ``the structure of our soul'' makes itself felt.} It has been objected
against this by several critics (in \textit{Revue philosophique}, 1904, and in \textit{American Journal of
Philosophy}, 1906) that the concept of continuity is not the same within the four problems. To this I answer
that in all domains there is only analogy, not identity, between the different applications or examples of
the same concept. Continuity appears in a different manner in the formal and in the real sciences, and
therefore also in epistemology on the one side, psychology and ethics on the other. And metaphysics stands
as a last attempt to carry through the concept of continuity, in that it lets a part or side of existence
represent the whole of existence. But there are in the end lacking means of thought to carry through the
connection. We stand at the boundary with the one compass-leg, but seek in vain a place for the other. No
rational analogy can be carried through here. And yet the type of thinking is the same. We always think in
predicates, but we cannot here find any subject to attach the predicates to.

When the concept of metaphysics is held fast in its strictest sense, it is to designate a seeking after an
absolute underlying ground that does not again point out beyond itself. Absoluteness must not be confused
with objectivity. For the psychologist soul-life has objective character, but spiritualism makes it
metaphysical by maintaining that in it we have the innermost essence of existence revealed. For the
physicist atoms and electrons have objective character, but the atomic doctrine first becomes metaphysical
when it (as in Democritus and others) contains the solution of all riddles.\footnote{Cf.\ on this last
point statements by Meyerson and Perrin in \textit{Bulletin de la société française de philosophie}, 22
January 1910.}

\section*{Excursus on the Term ``Metaphysics''}

As is well known, the term ``metaphysics'' arose from the circumstance that Aristotle's work on what he
called ``the first philosophy,'' the doctrine of the in-itself-existing, apart from its special forms, was in
a later collection of his writings placed next after (\textit{meta}) his ``Physics.'' The term occurs in a
single Neoplatonist and later in Arabic and medieval philosophers. Thomas Aquinas understood by metaphysics
a doctrine of the supernatural (\textit{res divinæ}) --- that is, of God and of the soul's immortal essence,
in so far as such a doctrine can be acquired by natural man. He distinguishes between it and positive
theology, which builds on revelations that have fallen to the lot of single chosen men. In more recent time
the word was often used of the doctrine of the spiritual, in opposition to natural science as the doctrine
of the material. But after psychology had been founded as an empirical science, metaphysics became a
doctrine of what lies beyond all possible experience at all, and the critique of the eighteenth century led
more and more to doubt of the possibility of such a science.\footnote{Cf.\ \textit{Vocabulaire technique et
critique de la philosophie} (Bulletin de la société française de philosophie, 1910).}

Besides the rising general tendency to lay weight on the science of experience, two thinkers have had great
influence on the fate of ``metaphysics'' --- Kant by his critique of the proofs of God's existence and of
the attempts at spiritualistic psychology, and Comte by his conception of ``metaphysics'' as an intermediate
link in the development of human cognition between theology and positive science.\footnote{On Comte's
conception of ``metaphysics'' see \textit{Den nyere Filosofis Historie}, 3rd ed., IV, p.\ 39\,f.} There is
here particular reason to go into Kant's use of the word, since he is not consistent in his usage.

Kant stated in one of his precritical writings that he had the misfortune to be in love with metaphysics, but
at the same time he ironized (in \textit{Träume eines Geistersehers}) over how easy it was at bottom to make
a science about the high matters metaphysics is to concern. The after-effect of his being in love still lays
itself before the day, as I have shown in my treatise on ``the continuity in Kant's philosophical
development'' (Vid.\ Selsk.\ Skr., 1893), in \textit{Kritik der reinen Vernunft}. He constantly maintained the
thought of the unity of nature, not as a mystical being lying behind the whole, but as an objective
correlate to the cognizing subject's urge and capacity to comprehend its observations into a rational whole.
Further, one can find an after-effect of his youthful love in the dogmatic manner in which he apprehended the
results of his epistemological investigations, in that he supposed he had furnished a proof of the validity of
the causal proposition, and finally also in his faith in having given a complete list of cognition's
fundamental concepts and fundamental propositions. As regards the very train of thought in \textit{Kritik der
reinen Vernunft}, it has recently\footnote{In Heimsoeth, \textit{Metaphysische Motive in der Ausbildung des
kritischen Idealismus} (Kant-Studien, 1924).} been shown that there lies a ``metaphysical'' foundation for
Kant's epistemology in the for him decisive opposition between passivity and activity --- an opposition that
led him to the unfortunate assumption of a ``Ding an sich'' as cause of our cognition, in so far as we
conduct ourselves passively. This demonstration strikes a point that was already drawn forth by the most
acute critique Kant got in his own time. Salomon Maimon namely showed that no point can be discovered in our
cognition where we conduct ourselves purely passively. Only approximately is there a purely passive
sense-sensation.\footnote{\textit{Den nyere Filosofis Historie}, 3rd ed., III, p.\ 122\,f.} The more recent
psychology confirms Maimon's assertion.

As regards Kant's own use of the word metaphysics, it has caused obscurity that he in the second edition of
\textit{Kritik der reinen Vernunft} calls the discussion (``Erörterung'') by which it is to be shown that
space and time are \textit{a~priori} forms of intuition ``metaphysical.'' In reality this discussion consists
in a psychological analysis that leads to singling out what lies in the sense-sensation in and for itself and
what is due to reflection (``the understanding''). This analysis then finds something that is neither
sensation nor reflection, and which is called ``form of intuition.'' According to Kant we in fact find space
and time as elements in our immediate sensing (besides the sense-sensation itself). In opposition to this
``metaphysical'' discussion Kant sets (in the second edition)\footnote{In the first edition the
epistemological result is given, after a psychological discussion of space and time (without heading), under
the heading ``Schlüsse aus den obigen Begriffen.''} a transcendental (that is, epistemological) discussion
that leads to seeing that space and time are necessary presuppositions for mathematical science. --- But in
the preface to the second edition Kant calls the whole demonstration of science's presuppositions that
\textit{Kritik der reinen Vernunft} gives ``the first part of metaphysics,'' which does not give a system, but
is only a ``treatise on method.'' As regards ``the second part of metaphysics,'' which was to treat of what
goes beyond the boundary of possible experience, Kant points out that there is now, after the first part has
shown the boundaries of science, free room for a faith to which ethical motives can lead. Kant's ``metaphysics,''
first part, thus corresponds to what we would call epistemology; its second part would belong under what we
call the philosophy of religion.

Unfortunately philosophical usage constantly presents wavering on this point as on other points. Thus Francis
Bradley calls his profound work \textit{Appearance and Reality} ``a metaphysical essay.'' In my eyes this
work is just as well an epistemological investigation as the section of \textit{Kritik der reinen Vernunft}
that Kant calls the doctrine of ``ideas'' (that is, of concluding concepts). And Bradley's work is then also
really a continuation of the named section of Kant's work.

\hnum{27} When an urge stirs to form representations that could give a last solution of the riddles of
cognition and of life, there must, according to what has been developed in the foregoing, be made a choice
between the different qualities or sides that existence empirically presents, and the chosen trait then stands
as characterizing existence in its innermost essence. If we hold to the most carried-through attempts in this
direction, two groups can be distinguished. The one holds to the fact that scientific cognition is possible to
the degree it is; the other lays the weight on a single one of the subjects existence presents for our
cognition.

In the first group it stands as an essential trait in existence's character that it is possible to reach a
comprehensive cognition whose presuppositions are again and again confirmed by new experiences and inferences.
Thereby existence acquires for us a rational character. If there were not unity and continuity in existence, it
could not come to meet our urge to cognition in the way it in fact does. If existence were a chaotic series of
difference --- that is, consisted of mutually equally different subjects that could be arranged in any way
whatever --- it would be irrational at its ground. But now there is an advancing work of thought that leads
through a series of stages to unity (identity) and whole (totality), and this is a witness of existence's
nature. Profound thinkers have interpreted this witness so that, as we think existence, existence thinks itself
in us. Our thought is itself a part of existence, and its progress despite all outer and inner hindrances bears
witness to the nature of the whole. --- Unity and whole can, as we have seen in the foregoing, step into
opposition to each other. From the standpoint of unity or identity the disparate and manifold that cannot be
reduced stands as a chaos. Parmenides held all attempts to order this chaos to be ``the opinions of mortals,''
and modern mathematicians and physicists call all representations of qualities and totalities
``anthropomorphisms,'' just as the theologians use this word for the psychological elements in the concept of
God, when in their opinion analogies with human psychology fail. But it has shown itself in the foregoing that
elements in the face of which the principle of identity is in and for itself impotent can yet join together
into wholes of different kinds.

Our cognition constantly works both in the direction of unity and in the direction of whole. It is a struggle
with chaos, and it wins its victory either by reducing differences to identity or by demonstrating real
connection between the different elements among themselves. Metaphysics becomes possible only by the fact that
the concept of unity and the concept of the whole are drawn out of the connection in which they occur in the
real work of cognition. And then two metaphysical systems can step into opposition to each other, in that the
one makes unity, the other the whole, into the absolute. --- In theology the opposition between God and the
world corresponds to the metaphysical opposition between unity and whole.\footnote{This train of thought partly
corrects the train of thought in \textit{Den menneskelige Tanke}, pp.\ 319--328, where the whole weight is
laid on the opposition between monism and pluralism.}

The other group of metaphysical attempts of thought lays all weight on a part of cognition's content. But
within this content there appear oppositions that have often stood as so enigmatic that one has supposed one
had to choose between their terms.

The opposition between spiritual and material has often pressed forward since Plato first drew it out. And the
problem was sharpened when the more recent natural science had asserted itself as a self-coherent world of
thought. The problem of spirit and matter posed itself when it seemed clear that the material could not be
derived from the spiritual, nor this from that. Nonetheless it became clear to clear thinking that the relation
between them was not contradictory, but that this unity is purely empirical and does not exclude other forms of
existence, which yet do not step forth in our experience. Spinoza saw this clearly: ``We neither sense nor
apprehend any single things other than bodies and phenomena of consciousness'' (\textit{Ethica} II, Axiom 5).
Perhaps it is then merely the limitation of our experience that makes us know only the two forms of existence.
And Spinoza supposed himself thus warranted to say that there are infinitely many forms of existence
(``attributes''), of which we in fact know only two.\footnote{See further \textit{Spinoza's Ethica. Analyse og
Karakteristik}, pp.\ 22--24.}

But since we have no possibility of deciding whether there really are other forms of existence, it is understood
that one for the most part goes other ways. One tries despite everything to derive the one form from the other,
or to let them stand in reciprocal action with each other, or to derive both from a third. All these ways one
has tried, but the history of thinking shows that none of them can be gone to the end. If we then, as with the
preceding problem, take our position within experience, it shows itself that the question really concerns the
relation between two sciences of experience, physiology and psychology. Physiology works according to the
natural-scientific method and meets, so long as this method is followed, nowhere a spiritual thing on its way.
And psychology explains the one spiritual problem by the other, as far as it is possible. But when the two ways
of investigation are compared, there appears a rational analogy between the two series of subjects, such that to
a subject within the one series there corresponds a definite subject within the other. There are then set the
two sciences tasks enough in pursuing this analogy as far as possible, and there is no ground for any kind of
metaphysics.\footnote{Cf.\ further \textit{Begrebet Analogi}, pp.\ 109--119.}

In a similar manner we must place ourselves in the face of another pair of opposites --- subsistence and
development (becoming, motion). Plato already, at a certain point of his development of thought, reached the
insight that the relation between rest and motion was not contradictory,\footnote{See \textit{Bemærkninger om
den platoniske Dialog Parmenides} (Vid.\ Selsk.\ Filosofiske Meddelelser, 1920), p.\ 56\,f.} without, however,
drawing all the consequences of this. It was first Leibniz who deepened this thought by apprehending rest as a
state in which moving forces hold each other in equilibrium. It will then --- when one passes over to
metaphysics --- become the question whether the tendency to equilibrium or the tendency to continued motion is
predominant in the world, and about this only continued experience will be able to instruct us. Here too we are
thus referred back, as with the other metaphysical problems, to the domain of experience and of the special
investigations. A fruitful treatment of the problems that have been considered ``metaphysical'' takes place only
when the posing of the problem acquires such a character as only the science of experience can give it.

\hnum{28} It has more and more shown itself impossible to find a concluding predicate for existence when one
tries to apprehend it as a whole. One has then supposed that what thought could not reach by its own power one
could find expression for by listening to life's own voice. Life has, after all, brought forth thought from its
womb! Schopenhauer already seized this expedient when, in opposition to speculative philosophy's logical
constructions, he appealed to ``\textit{der Wille zum Leben}'' [the will to life]. In a similar direction go
Nietzsche's slogan ``\textit{Wille zur Macht}'' [will to power] and Bergson's ``\textit{élan de vie}.'' In a
very fine and characteristic manner this tendency appears in Georg Simmel's treatise ``Vorformen der Idee'' (in
the journal \textit{Logos}, VI, 1906) and in his later treatise ``Der Konflikt der modernen Kultur.'' Besides
what has already earlier been mentioned concerning Simmel's train of thought, it shall here be named how he, from
the relation between psychology and epistemology, supposes himself to find a foundation for a new metaphysics.

Life, says Simmel, brings forth culture, and culture appears always in certain forms that correspond to given
conditions of the time. But life's unfolding continues its course and constantly leads to new forms. It is this
constant unfolding of life, the source of all forms, that from the first has led to science --- that is, to a
cognition that is not merely a means for human self-preservation, but shows itself to be a power that demands
recognition for the rules according to which it works. What life first brought forth for its own sake acquires,
by a kind of axis-rotation, independent significance, and life then becomes the servant of what it has itself
created. Such an axis-rotation takes place not merely when science comes to be as science, but also when art,
religion, and law come to be as independent spiritual powers. The ``axis-rotation'' itself shall, according to
Simmel, not be explicable by the psychological way, by a displacement of motives --- that is, by the fact that
what was first a means afterward becomes an end. For there happens by the ``axis-rotation,'' he supposes, a
complete liberation from the domain where there can be talk of means and ends. But although psychology and
epistemology are separated by an absolute boundary, there yet appear in psychological experience ``Vorformen,''
embryonic stages, for what afterward steps forth with ``overpsychological ideality.'' Thus practical insight is
an embryonic stage to scientific cognition. But it is life itself that now out of its womb brings forth logical
validity, objective cognition --- something that is more than life. Life thereby transcends itself. --- Now, here
as in other domains, continues Simmel, the tragic relation can set in that the forms, although they have proceeded
from life, no longer suffice, because life has advanced further. Hereby is understood the aversion to systematics
and to all forms whatever that is characteristic of our time. What formerly stood as norm, as objective truth, is
now regarded as something merely subjective, and a higher objectivity is sought. So it goes in the domain of art,
religion, and law, and so it goes also in the domain of thought.

It is the problem of the transition from psychology to epistemology that, according to Simmel, demands a
metaphysical solution. But his appeal to ``life'' has rather a mystical character. And there is here no ground for
mysticism. It is psychology's task to investigate the thinking that life led one to exercise, and which is itself
a kind of life, as far as it is possible. And the turning-point (the ``axis-rotation''), when thinking, so to
speak, takes the power from life, it also falls to psychology to describe and discuss as well as it is able. The
prescientific consciousness already involuntarily applies certain fundamental concepts (I call them the
fundamental categories), which science during its development more closely forms according to the demand of the
subjects. But this formation too psychology seeks, supported by the history of science, to understand. Here
difficulties and problems enough can remain, but they are not solved by a mystical appeal to life's creative
power. In any case there is continuity between the stage Simmel calls a thought's embryonic state and the one
where thought steps up with logical authority. One will constantly be able to find more intermediate links. And
now this logical authority is, according to Simmel, not always definitive. Its time can be past; the thought then
becomes what I call a dying category. It shows itself then that it was of subjective nature. But has it not been so
the whole time, even when it still stood in force? Is there any point in epistemology where it wholly lets go of
psychology, since it is, after all, constantly men who think even the most abstract and ideal thoughts, and since
all thoughts that are scientifically necessary must, after all, be psychologically possible?\footnote{A detailed
investigation of the relation between psychology and epistemology I have made in \textit{Totalitet som Kategori},
pp.\ 3--31.}

Even if it had been granted to Simmel to develop his above-cited indications further, he would yet surely never
have come into dogmatic metaphysics.\footnote{The metaphysical problem itself Simmel poses in a manner similar to
me: ``\textit{Aus den unübersehbar vielen Fäden, die das Netzwerk der Wirklichkeit ausmachen, und deren
Gesamtheit dem Philosophen sein Problem stellt, lässt ihn die Sonderart seines geistigen Typus einen einzelnen
ergreifen; ihn erklärt er für den, der das Ganze zusammenhält}'' [Out of the unsurveyably many threads that make
up the network of reality, and whose totality sets the philosopher his problem, the particular kind of his
spiritual type lets him seize a single one; this he declares to be the one that holds the whole together].
\textit{Hauptprobleme} (1910), p.\ 30.} He was convinced that a definitive form of satisfaction of the urge to
cognition will never be able to be found, and he adds: ``\textit{Es ist auch ein ganz philiströses Vorurteil, dass
alle Konflikte und Probleme dazu da sind, gelöst zu werden}'' [It is also a quite philistine prejudice that all
conflicts and problems are there to be solved] (Der Konflikt der modernen Kultur, p.\ 29).

\hnum{29} For the characterization of existence metaphysics is at bottom not at all necessary. Even if one
regards the single sciences' fundamental concepts and fundamental propositions as mere methods of work without
absolute validity, it yet casts light on the nature of existence that it can become more or less intelligible only
by the application of these concepts and propositions.

In his above-named work \textit{Appearance and Reality} (1893) Francis Bradley shows in a striking manner that no
subject and no kind of subject can stand as the expression of existence as a whole. What philosophy can reach is
not an idea that expresses the innermost essence of existence, but a measure for deciding which conception of
existence is the most perfect. This measure is drawn from the concept of experience. For complete experience there
would be required partly the richest fullness of subjects and qualities, partly the most intimate thought-connection
between these. By this measure we evaluate not merely the degree of perfection our experience has at a given time
reached, but also the existence that announces itself in this experience. The absolute existence would have an
unbounded content, and there would be a complete harmony between all elements in this content. Our real cognition is
constantly only an approximation in the direction of an experience in which such an existence would announce itself.

It is in the face of this clear and deep train of thought so much the more strange that Bradley can attribute no
significance at all, either to the methods and results of the science of spirit or to those of natural science, in
the attempt to make up thinking's position in the face of the problem of existence. But it has not hindered that a
renewed occupation with his works has convinced me that my thoughts go in the same direction as his, although I lack
the capacity to form so peculiar a union of mysticism and skepticism as the one that stamps his work.

\hnum{30} Already before Kant's critique overturned the assumption that science's fundamental thoughts expressed the
absolute essence of existence, Leibniz had, as mentioned above, maintained that all thoughts about the essence of
existence rested on analogy. In \textit{Den nyere Filosofis Historie} (3rd ed., II, p.\ 141) I say about this: ``The
way along which mythologies and earlier-set-up speculative systems, without being conscious of it, had come to the
assumption of the spiritual as ruling in existence --- this way Leibniz goes with full consciousness.'' But it cannot
yet be said that Leibniz saw the full reach of this thought. He did not see that there was more than one foundation
from which concluding analogies could be built. By the way of analogy a materialist metaphysics can be constructed,
not merely an idealist one. It depends on which side of experience one is most occupied with and has come to regard
to such a degree as the essential that one regards all other sides as derived in relation to it. The whole question of
metaphysics, also of metaphysics' special character, goes back to experience and to the side of experience that has
come to stand in the first line for thought, so that one by the way of analogy confidently lays it as the foundation
for one's view of the world.

Kant himself saw the great significance of analogies when one goes beyond every possible experience. In his discussion
of the concept of God which, according to his conception, it is warranted on ethical-religious motives to set up, he
declares (in all his three principal works) that when one takes all anthropomorphisms away from the representation of
God, only the mere word remains. That means that the analogies theology employs are not rational, but poetic. Just as
little as Leibniz did Kant see all the consequences of the fact that all concluding representations have a poetic
character.\footnote{Cf.\ \textit{Den nyere Filosofis Historie}, 3rd ed., III, pp.\ 95--97.}

When Kant supposed that there would, despite all critique, constantly stir in man an urge to metaphysics, it was
because he found in man's interior ``a hidden art,'' an urge and capacity to form the observations into a whole, a
capacity and an urge that in and for itself knows no boundaries, but continues its involuntary work even if experience
and science demonstrate definite boundaries for grounded cognition. Such involuntary constructions will yet not be of a
purely intellectual nature. Other spiritual needs too than the intellectual will play a role. Now Kant himself has, in
sharp opposition to the involuntary web of syntheses, set the strictly rational synthesis we call science, and it must
then come to a conflict between these two applications of the concept of synthesis. Psychology stands here otherwise
than epistemology. For the psychologist there is, as has rightly been said, no weed. He interests himself in all
thoughts, whether they can be used in science or not. Wild-growing representations too belong to his subjects, which he
by the help of experience and history seeks to understand. But epistemology investigates the conditions for the
representation's objective validity. The syntheses that are acknowledged are built on logical analysis of the subjects
and of the manner in which the subjects are connected. And the epistemological investigation will also direct itself
toward psychology's and history's attempts to explain the syntheses that are not built on logical analysis. --- There
will here not come to be a lack of problems. For conflicts between involuntarily developed assumptions and science's
critical demands will arise again and again, although under new forms.

In our days metaphysics (in the strict sense) is cultivated predominantly outside the circle of professional
philosophers. Wundt and Dilthey have already made this observation (see their treatises in ``Kultur der Gegenwart,''
I, 6). But there constantly stirs in many an urge to expansion (cf.\ \S\,19), to let thoughts one has once reached, and
moods one has once absorbed oneself in, run the line out. The philosopher can feel the same urge, but it is met in him
by the urge to a rational grounding of the representations he forms. The popular metaphysics presupposes a capacity to
look away from the difficulties experience and comparison can lead to. It takes no account of the difference between
intellectual and emotional thought-motives, or on the whole of the significance the division of labor in the spiritual
domain has acquired. In my treatise on Spinoza's Ethica (Vid.\ Selsk.\ Skr., 1918) I have, in the investigation of
Spinoza's three thought-motives, drawn attention to the fact that so intimate a union of different thought-motives as one
finds in the study of this thinker was possible only so long as the different problems had not yet been singled out from
one another. ---

The critique that must be directed against the attempts at metaphysics will also strike the ``philosophy of life'' that
supposes itself to have a shorter and easier way to truth than the one epistemology can acknowledge. It has rightly been
said of it that it solves the problems without posing them, and that it sets great words instead of the work of reaching
real insight.\footnote{Jonas Cohn in \textit{Philosophische Monatshefte}, 1925, pp.\ 13--15.} And yet such attempts can
have psychological and historical interest. We stand in the face of them, if they really have their ground in personal
experience and need, as we stand in the face of organisms and personalities. They can be acknowledged as naturally
grown wholes, but research reserves its analysis of them, and their worth will not least rest on whether, during their
growth, sufficient account has been taken of the cognition that under given conditions was at one's disposal. The
interest we can feel in the face of a faith borne by honest life-experience is akin to that felt in the face of a work of
art, but nonetheless the investigation of how it has arisen can be decisive for its value. Since metaphysics at bottom,
where it has lasting value, is a kind of poetry of life, the comparison with the work of art has a deeper significance.
We are all called to exercise the art of the view of life on the ground of experience, whether we exercise this art with
consciousness or not.

\chapter{Religion}
\label{ch:6}

\hnum{31} A view of the world, a metaphysics, acquires a religious character when it is a
practical-personal motive that leads to regarding a single subject within existence as the type of the
whole of existence. Conversely, every religion will have a more or less marked tendency to unfold a view
of the world. It thereby acquires a metaphysical character and often supposes itself able to solve the
riddles that science cannot solve. In religion as in metaphysics it is a single subject that is made into
the type of the whole. It may be the subsistence of the tribe, the condition for its security and its
nourishment. Or the motive may (in the religions that are called the higher) be the soul's need to know
its spiritual goods secure. The motive decides which subject comes to determine the religious view of the
world. Like all metaphysics, every religion gets its expression in symbols. The symbol corresponds to the
thought-motive and is an example that stands as the typical expression of the highest law in existence,
and about which the mind gathers itself in energetic concentration. The typical subject is here not merely
an object of contemplation, but is a source of consolation and breeds a model and a law for practical
striving. In the proper metaphysicians, on the other hand, abstraction and logical work of thought rule,
and therefore in them the opposition between part and whole, the chosen type and the rest of reality, will
more sharply make itself felt, so that the problem can stand more distinctly and a peculiar oscillation
between mysticism and skepticism appear. For a religious conception, on the other hand, a reconciliation
will have been reached, even if one constantly notices, and perhaps fears, ``the old enemy.''

Religious feeling always has a more or less complex character; quite simple it cannot be when it has
developed out of life-experience. It is due to the experience of hope and fear, of joy and sorrow, of
admiration and disgust, of exaltation and humiliation, of harmony and disharmony --- feelings called forth
during the course of life by the fate to which the highest valuable shows itself to be subject. Religious
feeling presupposes what one might call religious synthesis, a comprehending and comparing of value and
reality --- more exactly: of valuable reality and valueless or value-hostile reality. On the relation
between the two groups of experiences thus comprehended and compared, when religious feeling comes to be,
rests a religion's peculiarity --- for example, the character of the gods man honors and worships. As with
every total feeling,\footnote{Cf.\ on the concept of total feeling \textit{Den store Humor}, pp.\ 25--40.}
the synthesis can take two different forms: either fusion, so that the difference of the elements is not
noticed, or organization, so that one element as ground-value determines the other elements' influence.

That there is something valuable in the world even the most extreme pessimist cannot deny. His indignation
at existence's factual constitution presupposes that this constitution does not, after all, fill out all
his experience. He must stand at a place outside the bad world in order to be able to pass his disparaging
judgment on it. He bears witness, by his mere existence and his striving to judge the world and perhaps to
ground this judgment, that all is not of the evil. Goethe's Tasso counts it as the poet's advantage to be
able to utter his suffering, while other men fall silent in their pain. But in this the pessimist is right,
that the world of reality cannot be derived by the help of value-concepts, so that everything that subsists
should be form or means for something valuable. Plato posed all religion's great problem when he tried to
lay ``the idea of the good'' as the foundation for all existence. (Cf.\ above, \S\,1.) And yet the valuable
is not a foreign guest in the world that should suddenly and inexplicably emerge in single places. On
closer consideration value-concepts show themselves (as often mentioned in the foregoing) to hang closely
together with whole-concepts. So long as wholes form within nature, especially personal wholes, so long
will there be a foundation for the formation of value-concepts. What momentarily or constantly favors a
whole's subsistence acquires, from this whole's standpoint, value. A struggle between different wholes
(organic, personal, social) is a struggle between values, and here, as earlier mentioned, there are problems
enough. Of existence as a whole the value-concept cannot be applied, since it does not lie before us as a
given whole. Purely empirically the valuable stands as a stream that works its way forward and, despite
many hindrances, seeks to deepen and widen its bed. It is an image, but without images no view of the
world, whether it is religious or merely metaphysical, can unfold itself. Religious faith can be said to be
faith that that stream will constantly flow, and purer and freer than before. And since human work can work
along to secure and ease the stream's course, an ethical and a religious conception can here meet.

Since both the experiences and the personalities that make the experiences are different, the
image-formation that religious synthesis leads to will be able to take many different forms. There will
often, through the images employed, be traceable which personal types underlie. Thus at very different times
a living feeling of the resistance the valuable meets in reality has led to the use of analogous images: the
frog in a dry well in the Upanishads; the fish on dry land in St.\ Teresa and in Søren
Kierkegaard.\footnote{Cf.\ \textit{Religionsfilosofi}, 3rd ed., p.\ 73\,f.} A similar experience to the one
expressed by the image of the stream underlies a modern mystic's image of the plank that is rocked by the
waves, now gilded by the evening sun, now lying in darkness, and an older mystic's image of the iron that
becomes glowing in the fire, but otherwise lies as a black mass.\footnote{Cf.\ \textit{Oplevelse og Tydning}
[Experience and Interpretation], p.\ 64 and 97\,f.}

In the so-called positive religions free symbolizing can play only a subordinate role in relation to symbols
that are fixed once and for all and demanded to be apprehended as literal expressions of truth. In positive
religion grounded on tradition, personal life-experience is namely secondary. Certain teachings are handed
down, and the single individuals are in their life-experience to refind precisely what these teachings
present as serving the good of their souls. Peculiar personal life-experience and the image-formation
springing from it cannot indeed be kept away, even in the most authoritatively formed religion, since no
theology is able to utter thoughts that can acquire significance for all single individuals without being
individualized through personal life-experience and fantasy. But there will constantly anew arise a conflict
between individual experience and positive tradition. Life begins anew with every generation, and there
arises new need, new tasks, and new difficulties, for which there can be no remedy in the tradition. This
shows itself quite especially when a religion has been founded under cultural conditions very different from
the conditions of the time under which the single individual is to conduct his life. He cannot then find in
the tradition the models, the norms, or the consolation he needs. Even if he holds fast to the confession he
has once lived himself into, his real view of life will not be able to be covered by it. As already
mentioned in \textit{Begrebet Analogi} (p.\ 135), Stanley Hall has sought, by the help of the more recent
psychology, to give a picture of Jesus that should lead beyond the opposition between orthodoxy and
criticism and also show that the more recent time's difficulties and tasks are solved by Jesus standing as
humanity's ideal. But a real unity of character the Jesus does not, after all, acquire who is thus
constructed by combination of the old legends with the more recent time's culture. Stanley Hall is then also
clear that what constantly exercises and has exercised the greatest effect on men's minds is the Easter
cult, as it has been developed within the Church. But there it is not a constructed historical Jesus, but
the suffering God proclaimed by the Church on the foundation of the legend, that works. The opposition
between primitive Christianity and the more recent time's spiritual life is too great for an attempt to unite
both in a unity-picture to be able to succeed.

When such attempts are again and again made, it hangs together with what one has called our time's religious
unrest. Georg Simmel has, as mentioned above, characterized not merely the religious unrest, but the whole
unrest peculiar to our time, as dissatisfaction with the old cultural forms, without new ones being found.
In the religious domain he supposes that in very many men the religious faith in a beyond is now abandoned,
but without the will to religion (\textit{das religiöse Wollen}) therefore having fallen away.\footnote{\textit{Der
Konflikt der modernen Kultur}, p.\ 25.} When the will to religion is here said to subsist while the religious
content vanishes, what can this will then be other than an urge to get the better of life's oppositions, and
then especially the greatest of all, that between the valuable and the merely real? This urge then subsists,
while the capacity to satisfy it, which hitherto was supported by venerable tradition, cannot yet unfold
itself on its own. If such an unfolding is not possible, the urge too will in the end fall away. No
psychological or historical theory can guarantee its constant subsistence. But it is yet not excluded that
new religious geniuses can arise, in whom need and capacity work together in such a way that they can find
new spiritual positions in life's great struggle. Hitherto all great religious geniuses have come from the
Orient, but this has perhaps now played out its role. Europe and America live under the last rays of the
light that came from the East. But it would not be impossible that in the midst of the West's culture new
lights could arise. Should there come a spiritual movement of the kind indicated, it will hopefully not have
the shyness of problems and of independent reflection that Oriental religions have hitherto contributed much
to nourishing. Intellectual honesty will in earnest be acknowledged as a chief virtue that can least of all
be dispensed with in spiritual leaders.

Simmel's statement that in many the content of religion dies away while the religious urge constantly stirs is
illuminated by observations of Girgensohn.\footnote{Karl Girgensohn, \textit{Der seelische Aufbau des
religiösen Lebens} (1921), pp.\ 44--124.} By different ways he sought to get an insight into men's relation to
religion. The test-persons he reports on were non-churchly, but animated by a longing for ``ennoblement of
life'' --- for reaching a higher stage of life than the one on which they felt themselves to stand. They
declared that the word truth did not express what they longed for; what they wished were experiences that
could fill the mind.

Such an individual religious urge, devoid of content, as one has supposed oneself able to ascertain in our
time, is a modern phenomenon. Historically religion has essentially had a social character, in that it lays
itself before the day in forms of cult and in dogmas. One has even supposed that religion's proper content is
socially determined by the conditions for the tribe's subsistence and welfare.\footnote{The chief writing here
is constantly Durkheim, \textit{Les formes élémentaires de la vie religieuse} (1913). --- Cf.\ my review of
this work in \textit{Oplevelse og Tydning}.} Secure repose in handed-down cult and in the representations
conditioned by it is even today the first stage in the single individual's religious development. But a crisis
easily comes when the single individual, on his own, supported by his experiences and his reflection, is to
find his way in the world. Then the religious problem can come to be. It is due not to inventions and
accidents, but is an effect of the spiritual development that leads to advancing individualization. And this
problem bears the same stamp as the other problems the human mind, in virtue of its inner structure and its
experiences, comes to set itself. All problems come to be by a conflict between continuity and discontinuity,
and the religious problem is called forth by the question of the subsistence of value, although reality in so
many respects goes its own ways without regard to what men call values.

\hnum{32} Religion must not be judged merely by the factual forms that now have the greatest influence in the
world. Comparison with earlier forms casts light on the relation between the different elements that are
worked together in a positive religion at hand. In its most elementary forms religion has no system of dogmas
and no representations of gods, but it has a cult, through which magical influence is exercised on nature,
besides an inspiring effect being exercised on men's minds. In the central-Australian religion, which on the
basis of the description excellent observers have given is analyzed in an excellent manner by Durkheim, no
proper concept of God can be demonstrated. What makes it a religion is that it is governed by the great
opposition between holy and profane. Through the holy usages inherited from the fathers, man feels himself
raised above all individual striving and all care. Through the inherited cult, nature's power is also renewed.
Natural life, especially in so far as it serves to the good of men, would languish if there were not, through
the inherited cult's actions, brought to it new power. However striking these representations may now seem,
there is yet, as Lévy-Bruhl has pointed out,\footnote{\textit{Les fonctions mentales dans les sociétés
inférieures}, p.\ 295.} an analogy between them and the train of thought in higher religions' theistic
representations. To the periodic rejuvenation of nature by means of the observance of holy usages there
corresponds the necessity of God's constant intervention in nature in order for it to subsist. There is yet
this difference, that according to the theistic conception God subsists even if ``the world'' should no longer
subsist, while according to the primitive conception there is an unbroken reciprocal action between nature and
the magical powers the cult possesses.\footnote{Cf.\ \textit{Religionsfilosofi}, 3rd ed., pp.\ 42 and 163
(with note 115). --- From the scholastic side it has been objected against Spinoza's conception that it
assumes a relation of reciprocity between God and the world, while the right conception is that the world
indeed stands in relation to God, while God, who is eternal and unchangeable, does not stand in relation to
the world (\textit{Spinoza's Ethica}, p.\ 20\,f.).}

There are, as this example shows, analogies between the religions we call the higher and those we call the
lower. But analogy is not identity, and we can just as little without further ado interpret the lower
religions from representations drawn from the higher as conversely. Nonetheless the analogy here, as for
example in the relation between different sciences,\footnote{See \textit{Begrebet Analogi}, pp.\ 79--123.} can
in many respects be enlightening. ---

According to Durkheim the opposition between society and individual underlies the opposition between holy and
profane. Holy is what the tribe's welfare demands --- that is the primitive conception. But only so long as
individual and society are in solidarity, and so long as the handed-down representations of what the tribe's
welfare demands are not struck by criticism, only so long can this primitive conception subsist. For only so
long as the single individuals find in the social order, its traditions, and its usages sources of
inspiration, find their own urge to life and development satisfied, only so long can the tribe's good stand as
holy. The religious development then also leads from devotion to the factually subsisting society to faith in
a kingdom of God in which the perfect is to be realized.\footnote{The history of Jewish prophetism is here of
great interest. See \textit{Religionsfilosofi}, 3rd ed., p.\ 102\,f.}

The historical demonstration that the concept of God is not a primitive concept is of great interest from a
philosophy-of-religion standpoint. The first presupposition for the concept religion being applicable is that
there is something that stands as so valuable that man with tension follows its fate during the course of life.
Holy then becomes this valuable and everything that serves its subsistence. The religious judgments are
therefore secondary value-judgments. Primary value-judgments are such in which --- in the simplest form
through an exclamation --- it is established that there is something valuable. In the third line then comes the
urge to explain values and their subsistence by one or more beings' will and power.\footnote{See
\textit{Religionsfil.}, 3rd ed., pp.\ 61--65.} One can naturally narrow the concept religion so that it holds
only when the third presupposition is present. But it will from the side of research be an inexpedient
limitation. For all will be or become agreed that a concept of God has significance only when it is borne by
inward inspiration and reverence. But what qualities can awaken such inspiration and reverence? With this
question we return to the value-problem. What is it that makes God divine? It must be qualities attributed to
him on account of the valuable's coming-to-be and on account of its fate here in the world. Popularly one
apprehends God as cause of the world, which then stands as more or less valuable. Plato indicates a deeper
conception: that which makes God divine is that he can constantly live in the high ideas to which man only in
single, often ecstatic moments reaches up.\footnote{\textit{Platon's Bøger om Staten}, p.\ 133\,f.}

The main thing is that the word ``divine'' expresses a value-determination, and that the concept ``divine'' is
a predicate-concept. The very concept of God is, as we have seen in an earlier connection, a predicate-concept
before it becomes a subject-concept. First the qualities are discovered, thereafter the question can arise who
``has'' the qualities.\footnote{See further \textit{Religionsfilosofi}, 3rd ed., pp.\ 122--126. --- Cf.\
\textit{Oplevelse og Tydning}, pp.\ 88--91. --- Ludwig Feuerbach's philosophy-of-religion conception comes here
to its right (\textit{Den nyere Filosofis Historie}, 3rd ed., III, p.\ 267). --- Cf.\ in the most recent time
Poul Hofman, \textit{Das religiöse Erlebnis} (1925), pp.\ 60--62.} This consideration is supported from the
epistemological side by the fact that the concept ``essence'' (substance) is not a primary category, but is of
secondary significance. In the history of religion it shows itself time after time that it is the predicate
that is decisive. Perhaps the predicate is known beforehand, and one has waited with tension for the
possibility of applying it. Thus it was the primitive-Christian witness that the predicates ``Christ'' and
``Son of God'' could be applied to Jesus of Nazareth. In the Gospel of John XX, 31 we read: ``This is written,
that you should believe that Jesus is the Christ, the Son of God.'' About this predicate, its significance and
its validity, stands the strife between orthodoxy and historical-philosophical investigation.

\hnum{33} Different religious standpoints can be characterized through the religious synthesis on which they
rest --- that is, according to what has earlier been developed, according to the life-values that are known,
according to the conception of reality that is held, and according to the manner in which these two elements
are connected with each other.

A significant opposition arises from whether the world of values is apprehended as inferior in the face of the
world of reality, or whether the valuable stands as a power conquering despite all resistance. It is a
similar opposition to the one that in individual views of life appears between pessimism and optimism.
Pessimistic moods are for a great part due to the effect of contrast. The greater the demands our ideals set,
or the greater the expectations we come with, the greater the darkness we discover around us. And conversely,
the greater the resistance to our highest values, the brighter the mood toward life becomes when the
resistance is overcome. In his apology for Christianity Pascal seeks to show that Christianity, on its
supernatural foundation, lets life's great oppositions come to their right, and in his brilliant and profound
thoughts every serious view of life, even one that stands far from his supernatural presuppositions, can find
expression for the opposite moods that life awakens. Religions' cult has, even in their most elementary forms,
the power of letting men in great, ideal strokes live through the great oppositions. The Christian Easter
festival, with its opposition between Good Friday and Easter morning, is a mighty example of the psychological
hold that religion, by grasping and holding together life's most opposite experiences, has in the human mind.
It is a grand religious synthesis. On purely natural ground these opposite sides make themselves felt in what
I have called the great humor, which bravely takes up the great and the small, the exalted and the low, the
tragic and the comic as elements that experience, after all, presents, and which, perhaps only after hard
struggle, can also become elements in the ruling disposition. The great humor stands and falls with the
presupposition that what passes away is only single, empirical values, but that new sources of value will
again and again open themselves.\footnote{On the relation between humor and religion I must refer to
\textit{Den store Humor}, Ch.\ 5.}

The relation between value and reality is constantly experienced in new forms, as long as men carry on a
struggle among themselves and with the outer conditions. Generation after generation must here make its own
experiences. But under this struggle and these experiences the experiences of the past will constantly be of
significance. And the great religions are now precisely the expression of many generations' life-experiences,
interpreted according to the presuppositions that were at one's disposal. Whether a new generation can make
the same experiences and go in for the same interpretations, that is the religious problem that is set at
every time. But it is now not merely many generations' life-experiences that speak through religion; it is
also often different races' and peoples' experiences. For a great religion comes to be only by reciprocal
action between different races and different national religions. A striking example of this is Christianity, as
it, especially in the Pauline and the Johannine formation, spread in Europe and became dominant. Eminent
historians of religion have not been able to find this Christianity intelligible from the Jewish people and
its religion. Historical understanding is, according to them, here won only by assuming the influence of other
Oriental religions. From Jesus's personality as proceeded from this people, not all the predicates attributed
to it in the congregation are explained.

``The chief pieces of Christology stem,'' maintains Hermann Gunkel,\footnote{\textit{Zum
religionsgeschichtlichen Verständnis des Neuen Testamentes} (1903), p.\ 64.} ``not from the historical Jesus,
but have arisen independently of him and before him.'' Jesus made such an impression on his contemporaries and
posterity that one had to draw from other religious circles predicates to characterize him.

Such a historical conception cannot naturally be acknowledged from an orthodox standpoint. And philosophically
we must let it stand at its worth, leaving it to the historians of religion to test it more closely. But if
the named historical conception is correct, Christianity acquires, humanly seen, greater interest than after
the orthodox conception. For then the religion that victoriously penetrated into Europe is due to a sum of
experiences and interpretations that through long periods were lived through by several different races. The
purely human horizon is widened. When the liturgical part of the divine service nowadays wins in many greater
adherence than the single preachers' instruction, the explanation can lie in the fact that in the former one
hears many generations' more or less artistically formed witness, in the face of which a single speaker's
proclamation does not weigh. In the liturgy the old legends underlie, in which interpretations and events are
so mixed that no sharp boundaries can be drawn between them. But it is, after all, only through these legends
that Jesus has made his effect in the world. And I will here repeat what I have said in my
\textit{Religionsfilosofi} (3rd ed., p.\ 208): ``The most significant advances in the spiritual domain often
come from quarters from which one would least expect it. The new values very often come not from the world of
criticism and analysis, but from circles where one has with depth and inwardness lived in the old values.
History does not always go the straight ways.''

\hnum{34} That the psychological-historical standpoint for the religious problem is so strongly emphasized has
from the Catholic side\footnote{Goyau (see \textit{Religionsfilosofi}, 3rd ed., p.\ 202).} called forth the
question: Does religious truth come to be? Is it not something that exists outside the believers, stemming from
God? Or does it come to be in each single individual as a result of his life-experience? --- As I have already
stated in the cited place of my \textit{Religionsfilosofi}, the answer must consist in this, that in the
religious domain there can be talk of truth only in the sense that the representations one forms about
existence really correspond to what one has experienced about the relation between value and reality. Such an
experience must be made anew in every generation and by every single personality, and it is not necessary that
the result constantly become the same. Religious truth thus really comes to be. Religion's content is not
something that stands fast once and for all, and which one then goes and looks at. Within Catholicism such a
conception is not seldom. Imperial Chancellor Prince Hohenlohe, who was a Catholic, conceives in his memoirs
the relation between Catholics and Protestants such that, while the former can go and behold God's eternal
thoughts, which are revealed once and for all, the latter's religion must come to be anew in them. To this it
is yet to be remarked that the churchly Protestantism, which is afraid of its own consequences, does not
acknowledge the principle of personality, according to which the single personalities are the centers in which
life-experience and life-valuation, and thereby sincere religion, come to be. But if the principle of
personality breaks through as determinative for serious conception of life, it will be a witness of existence's
richness that it can be experienced and valued from so many starting-points. And through the conflict between
the different personal conceptions there will be able to arise new forms of conception of life, enriched by the
many sources, just as from a historical standpoint the great religions' coming-to-be is due to conflicts
between peoples and races. Hitherto it was essentially the reverse. The single individuals were to believe what
earlier generations believed, so that there was talk only of appropriation of traditions --- an appropriation
that became heresy as soon as it was undertaken with marked independence and on the basis of new experiences. In
stillness the principle of personality already now exercises its influence, an influence that lays itself before
the day not least in a growing capacity for mutual understanding, for finding and grasping the original and
peculiar, even in those whose view of life is another than our own. Such a capacity is the finest flower of all
culture.

\chapter{The Poetry of Life}
\label{ch:7}

\hnum{35} As a preparation for this last section of my investigations, a brief retrospect of the
preceding sections will be expedient.

When one arranges the sciences so that one begins with the most exact --- that is, the formal sciences,
which essentially build on identities --- and then advances to the real sciences, where qualities,
time-relations, and whole-forms play a more or less decisive role, one finds that, from standpoints drawn
merely from preceding links in this series, the following links come to stand as irrational or
``anthropomorphic.'' But as we have seen, endeavors make themselves felt to bring what in relation to
preceding links stood as irrational in under ``rational'' standpoints. If one cannot directly find
identities, one yet seeks to apply rational analogies. We have in this respect considered biology and
psychology, where it was especially the concept of the whole that made the application of the principle of
identity difficult. Identity and totality can step into sharp opposition to each other and thereby set the
great, in the end insoluble task of letting them both come perfectly to their right. Theology solves the
task by letting identity bring forth totality, or, in its language, letting God create the world. The
riddle of all riddles would undeniably be solved if we could form a concept of such a creating. But the
relation is in reality this, that our cognition works its way forward through attempts to apply the
principle of identity to the single parts of existence, in order thereby to approach a conception of the
existence accessible to us as a rational whole. And where we within the existence accessible to us find
factual whole-formations, we seek through analysis (thus likewise by the help of the principle of identity)
to understand the possibility of such whole-formations. This procedure is in the foregoing applied not
merely to the illumination of biology's and psychology's position, but also to the illumination of the
peculiar psychic totalities that appear as view of life, view of the world, and religion. It is these
domains that especially stand with the character of anthropomorphism, and a complete rationalization will
surely constantly be impossible here, however many attempts have been made to set view of life, view of the
world, and religion in a system. In the treatise on the concept of analogy it showed itself that beyond
the rational analogies that have significance for science, there are analogies due to the mind's urge to
find expression for the feeling of life aroused through the experiences made during the course of life ---
expression that can make them clearer to ourselves and make possible communication to others. Here there
takes place an image- or symbol-formation that becomes anthropomorphism's last refuge. Also in the
treatises on Totality and on Relation it showed itself that human spiritual life has not reached full
satisfaction of its urge there, where the rational application of these concepts must leave off.

Just as science has developed out of the natural life of consciousness through a purification and
precising of the (fundamental) categories that this already involuntarily applies in its reflection, so the
life of thought, when science can no longer lead it, returns to the natural ground from which it has
developed into science. But from this natural ground there can then, as experience shows, spring new, freer
forms of synthesis, and we have seen examples of such in view of life, in view of the world, and in
religion. From a philosophical standpoint such whole-formations stand as poetry of life. It is the urge of
self-preservation in single individuals and groups of individuals that here underlies and finds expression
in images and symbols. Faith in life and its values is, as has rightly been said by a Danish poet, the
greatest of all poets. And even if a psychological-historical analysis could throughout make such syntheses
explicable, their significance would not thereby vanish, when they really show themselves to satisfy what
personal truth and intellectual honesty demand.

\hnum{36} Poetry of life in the sense in which it is here taken stands in opposition to every conception of
life, religious or not, that fixes certain dogmatic assumptions as necessary. It seeks expression for life
as it stirs in us, conditioned as it is by what we have suffered, what we have enjoyed, and what we have
done. Its expressions are figurative, often formed involuntarily and at first hand, but often also
appropriated from art or from religious tradition. We take what we need, where we find it, when we cannot
ourselves produce it. Not merely great art, but also the great religions are products of the human spirit
and as such common property, of which everyone with right takes according to his need. Tegnér has stated
that poetry is the source of religion, and he calls (positive) religion ``circumcised poetry.'' Jacob
Burckhardt calls poetry the organ of religion, and Yrjö Hirn is, through his studies of medieval poetry and
art, led to the conception that in the representations of church doctrine there is traceable ``the effect of
an artistic creation that is no less significant for being unconscious and involuntary.'' According to Yrjö
Hirn, images and poetic works can often be understood only through ``the great, anonymous, and collective
poetry that one can read out of the whole churchly system of doctrine.''\footnote{In
\textit{Religionsfilosofi}, 3rd ed., p.\ 115 and note 57, I have already pointed to these statements of
Tegnér, Burckhardt, and Hirn.} In the churchly doctrinal concepts there are, as Tegnér's apt image
expresses, set definite boundaries for the involuntary unfolding of the poetry of life --- boundaries that
Yrjö Hirn's fine and deep research has broken through. Neither dogmatics nor criticism will in the long run
be able to hinder this unfolding. Involuntarily we compose ``a psalm of life.'' From the already mentioned
poem of Longfellow that bears this name, we see that in genuine poetry of life the single individual need not
feel himself isolated. We have the great spirits as models, who show us that it stands in our power to give
our life the stamp of the exalted (``Lives of great men all remind us, we can make our lives sublime''). And
we see at the same time that, however slight our deed has been, it can leave traces (``footprints on the
sands of time'') that will be able to guide others. It is not merely the great who become models; the small
too can become so, when they are faithful in their deed.

The poetry of life stands, to be sure, in a certain opposition to the poetry of art. It lies already in the
fact that it is due to an urge that is not always combined with artistic capacity. It will therefore often be
intelligible only to the single individual himself, and all others will, in the face of the manner in which he
finds expression for his urge to life and his life-experience, stand as profane ones who perhaps had best keep
away. There have in the foregoing been named examples of symbols that can have great significance for the
single individual, although for others they either mean nothing or cause offense. Such a relation has long
been known, for to the one religion the other's symbols are profane or unholy. Every man has, as Verner von
Heidenstam\footnote{``Hur lätt blir människornas kinder heta'' (\textit{Dikter}).} has said, his sanctuary
whose door is locked without anyone being able to find the key. The oil that burns in its lamp is secrets that
take place with himself.

But the poetry of art can itself step into the place of personal poetry of life and satisfy our urge through
its formed images of men and human destinies. He who himself, heavy and suffering, absorbs himself in
\textit{King Lear} finds there expression for fate's heavy course and hard hand. But when he comes to the end
and sees the few faithful ones stand around the corpses of the old king, shaken by his own passion and by
stern adversity, and his pure, bright daughter, his mood will acquire the character of inwardness and
exaltation, and he will sense what stirs in life's depths and on its heights. As with all tragedy there will
be brought about what one has called a ``purification'' of the feelings that our own striving calls forth in
us. And a feeling of gratitude will stir that a man has been able to find such an expression for life's
tragedy as we were not ourselves able to find.

In Danish literature the word poetry of life occurs first in Ingemann, whose youthful poetry of life
essentially bore the stamp of longing. Later Sibbern used the word of the wealth of lovely impressions that
daily life can bring us when we open eye and mind to them. Here it was not longing that was expressed, but joy
in what is nearest around us. And while Ingemann in his youth longed for the beyond, Sibbern maintained an
affirmation of life that takes up life as it once is.\footnote{See Vilhelm Andersen, \textit{Den danske
Literatur i det nittende Aarhundrede}, I, pp.\ 218 and 323.} --- It must yet be noted that later in Ingemann,
especially in ``Holger Danske,'' and quite particularly in ``Holger's Song to Life,'' manly tones came forth.
The powerful faith in the source of life, purified through resignation, that this last poem utters, and the
deep feeling of unity that in Paludan-Müller has found so beautiful an expression in ``Nattevagt'' [Night
Watch], are named by Troels Lund (\textit{Bakkehus og Solbjerg}) with right as excellent contributions to the
poetry of life from our literature's romantic period. But despite all it has rendered, the poetry of
Romanticism yet hovered high above the world of reality. One took the realities in nature and in history too
lightly. Its poetry of life is often only a reinterpretation of churchly representations, with a concealing of
what does not agree with these.

The poetry of life we need must have its root in free and deep life-experiences and spring from minds that are
well tried in the strife of life and of thoughts. It must not suppose itself able to solve problems at which
thought halts, nor be merely a paraphrase of handed-down representations. In bold flight it shall give great
symbols of that struggle of life, that conflict of light and darkness, that will continue as long as life
subsists.

Sometimes one has supposed that the time must come when all doubt was refuted, because science had solved all
riddles. But there will constantly arise new riddles, and it will be asked whether the solution of the old
riddles was correct. Science constantly works on a hypothetical foundation, from presuppositions that can
themselves again be made the foundation for investigation. And something similar holds for every view of life,
every faith. The richness of nature and of life will surely see to it that there will be new work of thought to
exercise, new struggles to stand, and new psalms of life to sing.

\end{document}
